% !TeX spellcheck = en_US
\documentclass[acmsmall,natbib,review,screen,anonymous,10pt]{acmart}
% TODO: FINALIZE
\settopmatter{printfolios=true,printccs=false,printacmref=false}

%\newcommand{\commentaire}{}
\newcommand{\MA}[1]{\ifdefined\commentaire{{{\color{green!60!black!55} 
\texttt{[#1]}}}}\fi}
\newcommand{\GM}[1]{\ifdefined\commentaire{{{\color[HTML]{5100FF} 
\texttt{[#1]}}}}\fi}
\newcommand{\GV}[1]{\ifdefined\commentaire{{{\color[HTML]{5100FF} 
\texttt{\small[#1]}}}}\fi}
\newcommand{\gb}[1]{\ifdefined\commentaire{{{\color{green!60!black!55} 
\texttt{[#1]}}}}\fi}
\newcommand{\bg}[1]{\ifdefined\commentaire{{{\color{green!60!black!55} 
\texttt{[#1]}}}}\fi}
\newcommand{\blue}[1]{\color[HTML]{0008FF} #1}


\newcommand*{\EHL}{\texttt{eHL}\xspace} %MA: abbreviation for expectation logic
\newcommand*{\HL}{\texttt{HL}\xspace}
\newcommand*{\PRHL}{\texttt{pRHL}\xspace}
\newcommand*{\ec}{\texttt{EasyCrypt}\xspace}
\newcommand{\easycrypt}{\ec}

\acmJournal{PACMPL}
\acmVolume{1}
\acmNumber{CONF} % CONF = POPL or ICFP or OOPSLA
\acmArticle{1}
\acmYear{2018}
\acmMonth{1}
\acmDOI{} % \acmDOI{10.1145/nnnnnnn.nnnnnnn}
\startPage{1}
\setcopyright{none}

\bibliographystyle{ACM-Reference-Format}
\citestyle{acmauthoryear}

\usepackage[utf8]{inputenc}
\usepackage{amsmath}
\usepackage{hyperref}
\usepackage{booktabs}
\usepackage{caption}
\usepackage{wrapfig}
\usepackage{subcaption}
\usepackage{multicol}
\usepackage{xspace}
\usepackage{xfrac}
\usepackage{tikz}
\usetikzlibrary{matrix,backgrounds,arrows.meta}
\usepackage{cleveref}
\usepackage{sbnf}
\usepackage{paper}
\usepackage{float}

\makeatletter
\let\@msm@th@eqref\eqref
\renewcommand{\eqref}[1]{%
  \begingroup
  \leavevmode
  \color{violet}%
  \hypersetup{linkbordercolor=[named]{violet}}%
  \@msm@th@eqref{#1}%
  \endgroup
}
\makeatother

\newcommand{\eqtag}[1]{\tag{{\color{violet}{#1}}}}


\captionsetup{subrefformat=parens}

\begin{document}

\newcommand{\todo}[2][red]{{\ttfamily\color{#1}[#2]}}
\newcommand{\ma}[1]{\todo[red]{MA: #1}}

\title{Lower bound reasoning with \textsf{eHL}}

\maketitle

\paragraph{The \EHL judgment.}
Let $\Mem$ denote the set of program memories and let
$\mathrm{Exp} = \Mem \to \Real^{+}$ be the set of
non-negative extended-real expectations.
For predicates $\pred, \predtwo : \Mem \to \Bool$,
expectations $\ex, \extwo \in \mathrm{Exp}$, and a program $\cmd$,
the \EHL judgment
\[
  \HT{\pred}{\ex}{\cmd}{\predtwo}{\extwo}
  \qquad {\rel} \in \{{\leq},=,{\geq}\}
\]

\paragraph{Validity.}
Let $\cmd_\mem$ denote the output distribution of $\cmd$ from memory $\mem$.
A judgment $\valid \HT{\pred}{\ex}{\cmd}{\predtwo}{\extwo}$ is \emph{valid}
if both of the following hold:
\begin{enumerate}
  \item $\forall\mem \in \Mem . \pred(\mem) \implies \supp(\cmd_\mem) \subseteq \predtwo$
  \item $\forall\mem \in \Mem . \pred(\mem) \implies \ex(\mem) \;\HTrel\; \E{\cmd,\mem}{[\predtwo]\cdot\extwo}$
\end{enumerate}
Condition~(1) requires that all output memories satisfy the postcondition
$\predtwo$. Condition~(2) requires that $\ex(\mem)$ bounds the expected
value of $\extwo$ restricted to outputs satisfying $\predtwo$.

\paragraph{Remark.}\label{rem:validity-indicator}
Since Condition~(1) guarantees $\supp(\cmd_\mem) \subseteq \predtwo$,
the indicator $[\predtwo]$ evaluates to $1$ on every output state of $\cmd$, and therefore
\[
  \E{\cmd,\mem}{[\predtwo]\cdot\extwo} \;=\; \E{\cmd,\mem}{\extwo}.
\]

\begin{figure}[H]
    \[
      \begin{array}{cc}
        \multicolumn{2}{c}{
        \Infer[hoare][skip]
        {\HT{\pred}{\ex}{\SKIP}{\pred}{\ex}}
        {}}
        \\[3mm]
        \multicolumn{2}{c}{
        \Infer[hoare][assign]
        {\HT{\pred[\vx \mapped \expr]}{\ex[\vx \mapped \expr]}{\vx <- \expr}{\pred}{\ex}}
        {}}
        \\[3mm]
        \multicolumn{2}{c}{
        \Infer[hoare][sample]
        {\HT{\forall \val \in \supp(\sexpr),\ \pred[\vx \mapped \val]}{\E{\sexpr}{\lambda \val.\, \ex[\vx \mapped \val]}}{\vx <* \sexpr}{{\pred}}{\ex}}
        {}}
        \\[3mm]
        \multicolumn{2}{c}{
        \Infer[hoare][seq]
        {\HT{\pred}{\ex}{\cmd;\cmdtwo}{\predtwo}{\extwo}}
        {\HT{\pred}{\ex}{\cmd}{\predthree}{\exthree} & \HT{\predthree}{\exthree}{\cmdtwo}{\predtwo}{\extwo}}}
        \\[3mm]
        \multicolumn{2}{c}{
        \Infer[hoare][if]
        {\HT{\pred}{\ex}{\ITE{\bexpr}{\cmd}{\cmdtwo}}{\predtwo}{\extwo}}
        {\HT{\bexpr\land\pred}{\ex}{\cmd}{\predtwo}{\extwo}& \HT{\neg\bexpr\land\pred}{\ex}{\cmdtwo}{\predtwo}{\extwo}}}
        \\[3mm]
        \multicolumn{2}{c}{
        \Infer[hoare][while]
        {\HT{\pred}{\lub\ex_{k}}{\WHILE\bexpr\DO\cmd}{\neg\bexpr\land\pred}{\extwo}}
        {\bexpr \land \pred \valid \ex_{0} \HTrel 0 
        & \forall k.\, \HT{\bexpr\land\pred}{\ex_{k+1}}{\cmd}{\pred}{\ex_{k}}
        & \neg\bexpr \land \pred \valid \ex_{k} \HTrel g}
        }
      \end{array}
    \]
  \caption{\EHL for lower and upper-bounds. Here ${\rel} \in \{{\leq},{\geq}\}$.}\label{fig:hoare:l-u}
\end{figure}

\paragraph{Notation.}
For a predicate $\pred : \Mem \to \Bool$ we use two indicator-style operators, one tailored to
each direction of reasoning:
\[
  [\pred](\mem) \;=\;
  \begin{cases} 1 & \text{if } \pred(\mem) \\ 0 & \text{otherwise} \end{cases}
  \qquad\text{(used for lower bounds)}
  \\[3mm]
\]
\[
  \langle \pred \rangle(\mem) \;=\;
  \begin{cases} 1 & \text{if } \pred(\mem) \\ \infty & \text{otherwise} \end{cases}
  \qquad\text{(used for upper bounds)}
  \\[3mm]
\]
% Given a memory-indexed expectation $r : \Mem \to \Real^{+}$,
% the expressions $[\pred] \cdot r$ and $\langle \pred \rangle \cdot r$ denote
% the corresponding pointwise products over $\Mem$:
% $[\pred] \cdot r$ agrees with $r$ on memories satisfying $\pred$ and is zero
% elsewhere, while $\langle \pred \rangle \cdot r$ agrees with $r$ on memories
% satisfying $\pred$ and is $\infty$ elsewhere.

\paragraph{The witness sequence $\{r_i\}_{i \in \mathbb{N}}$.}
The rule of \Cref{fig:hoare:while-lb} uses a sequence $\{r_i\}_{i \geq 0}$ of
non-negative expectations as a termination witness: intuitively, $r_i$ is a
(user-chosen) upper bound on how much of the post-expectation mass of $\ex$ is
still trapped inside the loop after $i$ body executions. Premise~2 ties $r_i$
to $\ex$ at $i=0$ by requiring $r_0 \geq \ex$; Premise~3 propagates the bound
across one iteration of the body; and Premise~4 forces the bound --- and thus
the trapped mass --- to vanish in the limit.


\begin{figure}[H]
  \small
    \[
      \Infer[hoare][while]
      {\HT[lower]{\pred}{\ex}{\WHILE\bexpr\DO\cmd}{\neg\bexpr\land\pred}{\ex}}
      {
        \HT[lower]{\bexpr\land\pred}{\ex}{\cmd}{\pred}{\ex}
        & \bexpr \land \pred \valid r_0 \geq \ex
        & \forall i \geq 0.\, \HT[upper]{\bexpr\land\pred}{r_{i+1}}{\cmd}{\top}{r_i}
        & \bexpr \land \pred \valid \lim_{i \to \infty} r_i = 0
      }
    \]
  \caption{The \textsc{While} rule for \textsf{eHL} lower bounds. The rule is
  specific to the lower-bound judgment: Premise~1
  is a lower-bound triple on the body, while Premise~3 is an
  \emph{upper}-bound triple used purely to bound the in-loop mass. The sequence $\{r_i\}_{i \geq 0}$
  consists of non-negative expectations $r_i : \mathrm{State} \to
  \mathbb{R}^+$.}\label{fig:hoare:while-lb}
\end{figure}


\newpage

\appendix

% ======================================================================
\section{Soundness and Completeness}\label{app:soundness-completeness}
% ======================================================================

\subsection{Foundations}\label{app:foundations}

This subsection makes explicit the typed semantic objects used by both the
soundness and the completeness arguments below.

\paragraph{Memories, predicates, and expectations.}
Let $\Mem$ be the set of program memories.  Predicates and expectations have
different codomains:
\begin{equation}
  P,Q:\Mem\to\Bool,
  \qquad
  f,g:\Mem\to[0,\infty].
  \label{eq:cook-semantic-types}
\end{equation}
All comparisons between expectations are pointwise; for example,
$f\leq g$ abbreviates $\forall m\in\Mem.\ f(m)\leq g(m)$.

\paragraph{Commands and distribution semantics.}
The command language considered below is
\begin{equation}
  C ::= \SKIP
  \mid \vx <- \expr
  \mid \vx <* \sexpr
  \mid C_1;C_2
  \mid \ITE{B}{C_1}{C_2}
  \mid \WHILE B\DO D.
  \label{eq:cook-command-syntax}
\end{equation}
Every command denotes a subdistribution-valued state transformer
\begin{equation}
  \llbracket C\rrbracket:
  \Mem\longrightarrow\mathcal D_{\leq1}(\Mem).
  \label{eq:cook-command-semantics-type}
\end{equation}
Subdistributions are used because non-termination may lose probability mass.
Two operations build every distribution used below, based on Avanzini
et al.\ [2024], Section~4. For $v\in\Mem$, write $\delta_v$ for
the \emph{point (Dirac) distribution} at $v$,
\begin{equation}
  \delta_v(m) \;\defsym\;
  \begin{cases} 1 & m=v, \\ 0 & \text{otherwise,}\end{cases}
\end{equation}
and define $\operatorname{return}(v)\in\mathcal D_{\leq1}(\Mem)$ by
\begin{equation}
  \operatorname{return}(v) \;\defsym\; \delta_v.
  \label{eq:cook-return-def}
\end{equation}
Given $\mu\in\mathcal D_{\leq1}(\Mem)$ and a family
$K:\Mem\to\mathcal D_{\leq1}(\Mem)$, the \emph{(monadic) sequencing}
operation $\operatorname{bind}(\mu,K)\in\mathcal D_{\leq1}(\Mem)$ mixes the
continuations $K(x)$ according to $\mu$:
\begin{equation}
  \operatorname{bind}(\mu,K)(m) \;\defsym\;
  \sum_{x\in\operatorname{supp}(\mu)} \mu(x)\cdot K(x)(m).
  \label{eq:cook-bind-def}
\end{equation}
With these two operations, the clauses needed most directly below include
\begin{align}
  \llbracket\SKIP\rrbracket(m)
    &=\operatorname{return}(m),
  \label{eq:cook-denote-skip}\\
  \llbracket\vx <- \expr\rrbracket(m)
    &=\operatorname{return}\bigl(m[\vx\mapped\expr(m)]\bigr),
  \label{eq:cook-denote-assign}\\
  \llbracket\vx <* \sexpr\rrbracket(m)
    &=\operatorname{bind}\bigl(\sexpr(m),
       \lambda v.\operatorname{return}(m[\vx\mapped v])\bigr),
  \label{eq:cook-denote-sample}\\
  \llbracket C_1;C_2\rrbracket(m)
    &=\operatorname{bind}\bigl(\llbracket C_1\rrbracket(m),
       \lambda m'.\llbracket C_2\rrbracket(m')\bigr).
  \label{eq:cook-denote-seq}\\
  \llbracket\ITE{B}{C_1}{C_2}\rrbracket(m)
    &=
      \begin{cases}
        \llbracket C_1\rrbracket(m),&B(m),\\
        \llbracket C_2\rrbracket(m),&\neg B(m).
      \end{cases}
  \label{eq:cook-denote-if}
\end{align}
For the loop, let $X:\Mem\to\mathcal D_{\leq1}(\Mem)$ and define the kernel
functional
\begin{equation}
  \mathcal F_{B,D}(X)(m)
  \defsym
  \begin{cases}
    \operatorname{bind}(\llbracket D\rrbracket(m),X),&B(m),\\
    \operatorname{return}(m),&\neg B(m).
  \end{cases} 
  \label{eq:cook-denote-loop-functional}
\end{equation}
Then
\begin{equation}
  \llbracket\WHILE B\DO D\rrbracket
  \defsym\operatorname{lfp}(\mathcal F_{B,D})
  \label{eq:cook-denote-loop}
\end{equation}
in the pointwise order on subdistribution kernels.

For a subdistribution $\mu$ and a non-negative function $h$, write
$\mathbb E_{x\sim\mu}[h(x)]\defsym\sum_{x\in\Mem}\mu(x)h(x)$ for its
expectation. Two elementary laws, used repeatedly below, follow directly
from the definitions of $\operatorname{return}$
(\ref{eq:cook-return-def}) and $\operatorname{bind}$
(\ref{eq:cook-bind-def}). The first is immediate, since
$\operatorname{return}(v)$ puts all its mass on $v$:
\begin{equation}
  \mathbb E_{x\sim\operatorname{return}(v)}[h(x)]
    =\sum_{x\in\Mem}\operatorname{return}(v)(x)\,h(x)
    =h(v).
  \label{eq:cook-expect-return}
\end{equation}
The second follows by unfolding $\operatorname{bind}$ and swapping the
order of summation:
\begin{align}
  \mathbb E_{y\sim\operatorname{bind}(\mu,K)}[h(y)]
    &=\sum_{y\in\Mem}\operatorname{bind}(\mu,K)(y)\,h(y)
    =\sum_{y\in\Mem}\sum_{x\in\Mem}\mu(x)\,K(x)(y)\,h(y)\\
    &=\sum_{x\in\Mem}\mu(x)\sum_{y\in\Mem}K(x)(y)\,h(y)
    =\mathbb E_{x\sim\mu}
       \left[\mathbb E_{y\sim K(x)}[h(y)]\right].
  \label{eq:cook-expect-bind}
\end{align}
The second equation is the law of iterated expectation for sequential
composition; all sums involved are of non-negative terms, so the
rearrangement is justified regardless of convergence.

\paragraph{The weakest-precondition operator.}
A single operator $\WP{C}{\cdot}$ is used throughout, applied either to an
expectation or to a predicate; which of the two defining equations applies
is determined by the type of the argument --- no separate notation is
introduced for the two cases:
\begin{align}
  \WP{C}{g}(m)
    &\defsym
      \mathbb E_{m'\sim\llbracket C\rrbracket(m)}[g(m')],
    &&g:\Mem\to[0,\infty],
  \label{eq:cook-wp-definition}\\
  \WP{C}{Q}(m)
    &\defsym
      \forall m'\in\supp(\llbracket C\rrbracket(m)).\ Q(m'),
    &&Q:\Mem\to\Bool.
  \label{eq:cook-gwp-definition}
\end{align}
Applied to an expectation, $\WP{C}{g}$ is quantitative: it returns the exact
expected value of $g$ after running $C$. Applied to a predicate, $\WP{C}{Q}$
is the ordinary (liberal) weakest precondition, expressing that every
reachable output state satisfies $Q$.

The return and bind laws give the corresponding equations for the predicate
case,
\begin{equation}
  \WP{\SKIP}{Q}=Q,
  \qquad
  \WP{C_1;C_2}{Q}
  =\WP{C_1}{\WP{C_2}{Q}}.
  \label{eq:cook-gwp-basic-laws}
\end{equation}
The identical-looking equations for the expectation case are
\eqref{eq:cook-wp-skip} and \eqref{eq:cook-wp-seq} below; both pairs are
instances of the same return/bind laws, one for each type of argument.

\paragraph{Semantic validity versus syntactic derivability.}
For $\bowtie\in\{\leq,=,\geq\}$, semantic validity of a guarded judgment is
\begin{equation}
  \valid\HT{P}{f}{C}{Q}{g}
  \quad\Longleftrightarrow\quad
  \forall m\in\Mem.\ P(m)\Longrightarrow{}
  \bigl(\WP{C}{Q}(m) \land f(m)\bowtie\WP{C}{g}(m)\bigr).
  \label{eq:cook-validity-typed}
\end{equation}
The first conjunct is qualitative: every output state in the support
satisfies $Q$. The second is quantitative: $f$ has the required relation to
the exact expected value of $g$. The two conjuncts are not independent: since
the first conjunct forces $[Q]=1$ on the whole output support, it gives
$\mathbb E[[Q]\cdot g]=\mathbb E[g]$, so \eqref{eq:cook-validity-typed} is
equivalent to the validity definition used at the start of the paper, which
weights $g$ by $[Q]$ explicitly instead of stating support-inclusion as a
separate conjunct. Throughout the rest of this section, $\vDash$ denotes
semantic truth and $\vdash$ denotes derivability from the displayed proof
rules.

\subsection{Soundness of the While Rule}\label{app:while-soundness}

\paragraph{Goal.}
We want to show that the \textsc{While-lb} rule of \Cref{fig:hoare:while-lb}
is sound, i.e.\ that its conclusion
$\HT[lower]{\pred}{\ex}{\WHILE\bexpr\DO\cmd}{\neg\bexpr\land\pred}{\ex}$ is
valid. Its postcondition is $\neg\bexpr\land\pred$, so by validity
Condition~(2) the indicator $[\pred\land\neg\bexpr]$ below is exactly the
postcondition indicator applied to $\ex$; we keep it explicit because it is
manipulated directly throughout the proof:
\[
  [\pred] \cdot \ex
  \leq
  \WP{\WHILE\bexpr\DO\cmd}{[\pred \wedge \neg\bexpr] \cdot \ex}
\]

\paragraph{Definitions.}
Define the family of approximating programs:
\[
  W_n
  \;\defsym\;
  \bigl(\kwd{if}\ \bexpr\ \kwd{then}\ \cmd\bigr)^n
  \mathbin{;}
  \kwd{if}\ \bexpr\ \kwd{then}\ \ABORT
\]
where $W_n$ allows at most $n$ iterations; if the loop would continue beyond $n$,
it aborts. Define:
\[
  I_n \;\defsym\; \WP{W_n}{[\pred \wedge \neg\bexpr] \cdot \ex}
\]

% ----------------------------------------------------------------------
\begin{lemma}[Recursive characterisation of $I_n$]\label{lem:In-recursion}
\[
  I_0 \;=\; [\pred \wedge \neg\bexpr] \cdot \ex
  \qquad\text{and}\qquad
  I_{n+1}
  \;=\; [\bexpr] \cdot \WP{\cmd}{I_n}
        + [\pred \wedge \neg\bexpr] \cdot \ex
\]
\end{lemma}

\begin{proof}
By induction on $n$.

\medskip\noindent\textbf{Base case $n=0$.}
\begin{align*}
  I_0
  &= \WP{\kwd{if}\ \bexpr\ \kwd{then}\ \ABORT}{[\pred \wedge \neg\bexpr] \cdot \ex} \\
  &= [\bexpr] \cdot \WP{\ABORT}{[\pred \wedge \neg\bexpr] \cdot \ex}
     + [\neg\bexpr] \cdot [\pred \wedge \neg\bexpr] \cdot \ex \\
  &= [\bexpr] \cdot 0 + [\pred \wedge \neg\bexpr] \cdot \ex
   \;=\; [\pred \wedge \neg\bexpr] \cdot \ex
\end{align*}
where we used $\WP{\ABORT}{\cdot} = 0$ and
$[\neg\bexpr] \cdot [\pred \wedge \neg\bexpr] = [\pred \wedge \neg\bexpr]$.

\medskip\noindent\textbf{Inductive step $n \to n+1$.}
Since $(\kwd{if}\ \bexpr\ \kwd{then}\ \cmd)^{n+1}
 = \kwd{if}\ \bexpr\ \kwd{then}\ \cmd
   \mathbin{;} (\kwd{if}\ \bexpr\ \kwd{then}\ \cmd)^n$, we have
$W_{n+1} = (\kwd{if}\ \bexpr\ \kwd{then}\ \cmd) \mathbin{;} W_n$, so:
\begin{align*}
  I_{n+1}
  &= \WP{\kwd{if}\ \bexpr\ \kwd{then}\ \cmd}{\WP{W_n}{[\pred \wedge \neg\bexpr] \cdot \ex}} \\
  &= \WP{\kwd{if}\ \bexpr\ \kwd{then}\ \cmd}{I_n} \\
  &= [\bexpr] \cdot \WP{\cmd}{I_n} + [\neg\bexpr] \cdot I_n
\end{align*}
It remains to identify $[\neg\bexpr] \cdot I_n$. When $\neg\bexpr$ already
holds at the starting memory, none of the guard checks inside $W_n$ ever
fire, so the memory is left untouched and $I_n$ reduces directly to
$[\pred \wedge \neg\bexpr] \cdot \ex$:
\[
  [\neg\bexpr] \cdot I_n
  \;=\; [\pred \wedge \neg\bexpr] \cdot \ex.
\]
Therefore:
\[
  I_{n+1}
  \;=\; [\bexpr] \cdot \WP{\cmd}{I_n} + [\pred \wedge \neg\bexpr] \cdot \ex
\]
\end{proof}

% ----------------------------------------------------------------------

\begin{lemma}[Key inequality]\label{lem:key-ineq}
Assuming premises \textnormal{(a)}, \textnormal{(b)}, \textnormal{(c)},
for all $n \geq 0$:
\[
  [\pred] \cdot I_n + [\pred] \cdot r_n \;\geq\; [\pred] \cdot \ex
\]
\end{lemma}

\begin{proof}
By induction on $n$.

\medskip\noindent\textbf{Base case $n = 0$.}
\[
  [\pred] \cdot I_0 + [\pred] \cdot r_0
  \;=\; [\pred \wedge \neg\bexpr] \cdot \ex + [\pred] \cdot r_0
\]
By premise~(b), $[\pred] \cdot r_0 \geq [\pred \wedge \bexpr] \cdot \ex$, so:
\[
  \;\geq\; [\pred \wedge \neg\bexpr] \cdot \ex + [\pred \wedge \bexpr] \cdot \ex
  \;=\; [\pred] \cdot \ex
\]

\medskip\noindent\textbf{Inductive step $n \to n+1$.}
Expanding $I_{n+1}$ by \Cref{lem:In-recursion} and multiplying by $[\pred]$:
\begin{align*}
  [\pred] \cdot I_{n+1} + [\pred] \cdot r_{n+1}
  &= [\pred \wedge \bexpr] \cdot \WP{\cmd}{I_n}
     + [\pred \wedge \neg\bexpr] \cdot \ex
     + [\pred] \cdot r_{n+1}
\end{align*}
Premise~(c) gives $[\pred \wedge \bexpr] \cdot r_{n+1} \geq \WP{\cmd}{r_n}$. Multiplying by
$[\pred]$ (using $[\pred]\cdot[\pred\wedge\bexpr] = [\pred\wedge\bexpr]$):
\[
  [\pred \wedge \bexpr] \cdot r_{n+1} \;\geq\; [\pred] \cdot \WP{\cmd}{r_n}
\]
Since $r_n \geq [\pred]\cdot r_n$ pointwise ($r_n$ is non-negative), monotonicity of $\WP{\cmd}{\cdot}$
gives $\WP{\cmd}{r_n} \geq \WP{\cmd}{[\pred]\cdot r_n}$, hence
$[\pred]\cdot\WP{\cmd}{r_n} \geq [\pred\wedge\bexpr]\cdot\WP{\cmd}{[\pred]\cdot r_n}$. Combining, and
extending from $[\pred\wedge\bexpr]$ to $[\pred]$ using non-negativity of $r_{n+1}$ on the
$[\pred\wedge\neg\bexpr]$ part:
\[
  [\pred] \cdot r_{n+1} \;\geq\; [\pred \wedge \bexpr] \cdot \WP{\cmd}{[\pred] \cdot r_n}
\]
Substituting into the expansion of $[\pred]\cdot I_{n+1}+[\pred]\cdot r_{n+1}$ above:
\begin{align*}
  [\pred] \cdot I_{n+1} + [\pred] \cdot r_{n+1}
  &\geq [\pred \wedge \bexpr] \cdot \WP{\cmd}{I_n}
        + [\pred \wedge \bexpr] \cdot \WP{\cmd}{[\pred] \cdot r_n}
        + [\pred \wedge \neg\bexpr] \cdot \ex
\end{align*}
By linearity of $\WP{\cmd}{\cdot}$:
\begin{align*}
  &= [\pred \wedge \bexpr] \cdot \WP{\cmd}{I_n + [\pred] \cdot r_n}
     + [\pred \wedge \neg\bexpr] \cdot \ex
\end{align*}
Since $I_n \geq [\pred]\cdot I_n$ pointwise ($I_n$ is non-negative), monotonicity of $\WP{\cmd}{\cdot}$
gives:
\[
  [\pred \wedge \bexpr] \cdot \WP{\cmd}{I_n + [\pred] \cdot r_n}
  \;\geq\; [\pred \wedge \bexpr] \cdot \WP{\cmd}{[\pred]\cdot I_n + [\pred] \cdot r_n}
\]
By the induction hypothesis
$[\pred] \cdot I_n + [\pred] \cdot r_n \geq [\pred] \cdot \ex$
and monotonicity of $\WP{\cmd}{\cdot}$:
\begin{align*}
  &\geq [\pred \wedge \bexpr] \cdot \WP{\cmd}{[\pred] \cdot \ex}
        + [\pred \wedge \neg\bexpr] \cdot \ex
\end{align*}
By premise~(a), validity Condition~(2) gives
$[\pred \wedge \bexpr] \cdot \WP{\cmd}{[\pred] \cdot \ex} \geq [\pred \wedge \bexpr] \cdot \ex$:
\[
  \;\geq\; [\pred \wedge \bexpr] \cdot \ex + [\pred \wedge \neg\bexpr] \cdot \ex
  \;=\; [\pred] \cdot \ex
\]
\end{proof}

% ----------------------------------------------------------------------
\paragraph{Main result.}
From \Cref{lem:key-ineq}, for all $n \geq 0$:
\[
  [\pred] \cdot I_n + [\pred] \cdot r_n \;\geq\; [\pred] \cdot \ex
\]
Taking the limit:
\[
  \lim_n ([\pred] \cdot I_n) + \lim_n ([\pred] \cdot r_n) \;\geq\; [\pred] \cdot \ex
\]
By condition~(d), $\lim_n r_n = 0$, so $\lim_n ([\pred] \cdot r_n) = 0$:
\[
  \lim_n ([\pred] \cdot I_n) \;\geq\; [\pred] \cdot \ex
\]
By definition and \Cref{lem:In-recursion},
$\lim_n I_n = \WP{\WHILE\bexpr\DO\cmd}{[\pred \wedge \neg\bexpr] \cdot \ex}$,
giving:
\[
  \WP{\WHILE\bexpr\DO\cmd}{[\pred \wedge \neg\bexpr] \cdot \ex}
  \;\geq\; [\pred] \cdot \ex
\]
which is exactly the desired conclusion. \qed

\subsection{Completeness via Exact Pre-Expectations}\label{app:completeness-cook}

This section proves completeness for guarded judgments with arbitrary
qualitative pre- and postconditions.  Let
$P,Q:\Mem\to\Bool$ and $f,g:\Mem\to[0,\infty]$.  For every command $C$ and
post-assertion $(Q\mid g)$, structural induction on $C$ derives the judgment
whose two precomponents are semantically exact: the qualitative component is
$\WP{C}{Q}$ and the quantitative component is $\WP{C}{g}$.  A final use of
\textsc{Consequence} then derives every semantically valid judgment
$\HT{P}{f}{C}{Q}{g}$.

Throughout,
\[
  \bowtie\ \in\ \{\,\leq,=,\geq\,\}
\]
denotes the lower, equality, or upper judgment relation. The induction is
carried out once, for an arbitrary $\bowtie$.

\subsubsection{Exactly what is proved}

For every command $C$, predicates $P,Q$, and expectations $f,g$, we prove
\begin{equation}
  \valid\HT{P}{f}{C}{Q}{g}
  \quad\Longrightarrow\quad
  \vdash\HT{P}{f}{C}{Q}{g}.
  \label{eq:cook-target}
\end{equation}
By the typed validity definition \eqref{eq:cook-validity-typed}, the premise of
\eqref{eq:cook-target} means
\begin{equation}
  \forall m\in\Mem.\quad
  P(m)\Longrightarrow
  \bigl(\WP{C}{Q}(m)\land
        f(m)\bowtie\WP{C}{g}(m)\bigr),
  \label{eq:cook-general-validity}
\end{equation}
or, equivalently,
\begin{equation}
  P\models \WP{C}{Q}\land f\bowtie\WP{C}{g}.
  \label{eq:cook-general-validity-entailment}
\end{equation}
The proof proceeds through the following exact-pre judgment:
\begin{equation}
  \vdash
  \HT{\WP{C}{Q}}{\WP{C}{g}}{C}{Q}{g}
  \label{eq:cook-general-exact-pre-target}
\end{equation}
Here $\WP{C}{Q}$ is the exact qualitative precondition for $Q$, and
$\WP{C}{g}$ is the exact quantitative pre-expectation of $g$.  Once
\eqref{eq:cook-general-exact-pre-target} has been derived, the desired judgment
follows by \textsc{Consequence}:
\begin{equation}
\frac{
  \begin{gathered}
    P\models\WP{C}{Q}\land f\bowtie\WP{C}{g}
    \\[1mm]
    \vdash\HT{\WP{C}{Q}}{\WP{C}{g}}{C}{Q}{g}
    \\[1mm]
    Q\models Q\land g\bowtie g
  \end{gathered}
}{
  \vdash\HT{P}{f}{C}{Q}{g}
}\;\textsc{Consequence}.
\label{eq:cook-general-final-consequence-outline}
\end{equation}
The first premise is \eqref{eq:cook-general-validity-entailment}, the second is
the exact-pre judgment, and the third holds by reflexivity.  It therefore
suffices to prove \eqref{eq:cook-general-exact-pre-target} by structural
induction on $C$.

\subsubsection{The rules and semantic equations used}

The \textsc{Skip}, \textsc{Assign}, \textsc{Sample}, \textsc{Seq}, and
\textsc{If} rules are those in \Cref{fig:hoare:l-u}.  They are parameterised
by $\bowtie$ and therefore cover lower, equality, and upper judgments at once.

\paragraph{The guarded consequence rule.}
Following the actual guarded consequence rule and the notation in the rest of
this note, we use
\begin{equation}
\frac{
  P\models P'\land f\bowtie f'
  \qquad
  \vdash\HT{P'}{f'}{C}{Q'}{g'}
  \qquad
  Q'\models Q\land g'\bowtie g
}{
  \vdash\HT{P}{f}{C}{Q}{g}
}\;\textsc{Consequence}.
\label{eq:cook-consequence}
\end{equation}
For $\bowtie={\leq}$ this is the lower rule, for $\bowtie={=}$ the equality
rule, and for $\bowtie={\geq}$ the upper rule.  The first comparison is
required on the new pre-guard $P$, and the last comparison is required on the
source post-guard $Q'$.  Every use of \textsc{Consequence} below is displayed
in this three-premise form with all entries already instantiated.

\paragraph{Semantic equations.}
The expectation semantics gives
\begin{align}
  \WP{\SKIP}{g} &= g,
  \label{eq:cook-wp-skip}\\
  \WP{\vx <- \expr}{g} &= g[\vx\mapped\expr],
  \label{eq:cook-wp-assign}\\
  \WP{\vx <* \sexpr}{g}
    &= \E{\sexpr}{\lambda\val.\,g[\vx\mapped\val]},
  \label{eq:cook-wp-sample}\\
  \WP{C_1;C_2}{g}
    &= \WP{C_1}{\WP{C_2}{g}},
  \label{eq:cook-wp-seq}\\
  \WP{\ITE{B}{C_1}{C_2}}{g}
    &= [B]\cdot\WP{C_1}{g}
       +[\neg B]\cdot\WP{C_2}{g}.
  \label{eq:cook-wp-if}
\end{align}
These equations follow from the denotational clauses above rather than being
additional proof rules: the first four follow from the return and bind laws,
and the conditional equation follows by splitting on $B(m)$.

For $L=\WHILE B\DO D$, define
\begin{equation}
  \Phi_g(X)\defsym[B]\cdot\WP{D}{X}+[\neg B]\cdot g.
  \label{eq:cook-loop-functional}
\end{equation}
The loop semantics used in this note is
\begin{equation}
  \WP{L}{g}=\operatorname{lfp}(\Phi_g)
  =\sup_{n\geq0}\Phi_g^n(0).
  \label{eq:cook-loop-lfp}
\end{equation}

We assume that the exact preconditions, pre-expectations, and loop residuals
used below are expressible, and that the background assertion theory proves
the required semantic equalities and inequalities.  The loop case additionally
uses monotonicity, Scott continuity, and additivity of $\mathsf{wp}$.  Its exact
pre-expectation is assumed finite-valued so that residual subtraction never
forms $\infty-\infty$.

\subsubsection{The exact-pre lemma}

\begin{lemma}[Exact-pre judgments]
\label{lem:exact-pre-derivable}
For every command $C$, every postcondition $Q$, every post-expectation $g$ in
the expressive, finitary fragment, and each
$\bowtie\in\{\leq,=,\geq\}$,
\begin{equation}
  \vdash\HT{\WP{C}{Q}}{\WP{C}{g}}{C}{Q}{g}.
  \label{eq:cook-exact-bowtie}
\end{equation}
\end{lemma}

\begin{proof}
The proof is by structural induction on $C$, with $Q$, $g$, and $\bowtie$
universally quantified.

\paragraph{Case: skip.}

The typed weakest-precondition equations are
\[
  \WP{\SKIP}{Q}=Q,
  \qquad
  \WP{\SKIP}{g}=g.
\]
Hence \textsc{Skip} gives
\[
\frac{}{\vdash\HT{Q}{g}{\SKIP}{Q}{g}}\;\textsc{Skip},
\]
which is \eqref{eq:cook-exact-bowtie} after rewriting both precomponents.

\paragraph{Case: assignment.}

The typed assignment equations are
\[
  \WP{\vx <- \expr}{Q}=Q[\vx\mapped\expr],
  \qquad
  \WP{\vx <- \expr}{g}=g[\vx\mapped\expr].
\]
The required instance of \textsc{Assign} is
\[
\frac{}{
  \vdash\HT{Q[\vx\mapped\expr]}{g[\vx\mapped\expr]}
    {\vx <- \expr}{Q}{g}
}\;\textsc{Assign}.
\]
Rewriting its two precomponents proves \eqref{eq:cook-exact-bowtie}.

\paragraph{Case: sampling.}

The typed sampling equations are
\[
  \WP{\vx <* \sexpr}{Q}
  =\forall\val\in\supp(\sexpr).\ Q[\vx\mapped\val],
  \qquad
  \WP{\vx <* \sexpr}{g}
  =\E{\sexpr}{\lambda\val.\,g[\vx\mapped\val]}.
\]
Thus \textsc{Sample} gives
\[
\frac{}{
  \vdash\HT{\forall\val\in\supp(\sexpr).\ Q[\vx\mapped\val]}
    {\E{\sexpr}{\lambda\val.\,g[\vx\mapped\val]}}
    {\vx <* \sexpr}{Q}{g}
}\;\textsc{Sample}.
\]
Rewriting its two precomponents proves \eqref{eq:cook-exact-bowtie}.

\paragraph{Case: sequential composition.}

Define
\[
  R\defsym\WP{C_2}{Q},
  \qquad
  h\defsym\WP{C_2}{g}.
\]
The two induction hypotheses are
\[
  \vdash\HT{R}{h}{C_2}{Q}{g},
  \qquad
  \vdash\HT{\WP{C_1}{R}}{\WP{C_1}{h}}{C_1}{R}{h}.
\]
They are exactly the two premises of \textsc{Seq}:
\begin{equation}
\frac{
  \vdash\HT{\WP{C_1}{R}}{\WP{C_1}{h}}{C_1}{R}{h}
  \qquad
  \vdash\HT{R}{h}{C_2}{Q}{g}
}{
  \vdash\HT{\WP{C_1}{R}}{\WP{C_1}{h}}{C_1;C_2}{Q}{g}
}\;\textsc{Seq}.
\label{eq:cook-seq-derivation}
\end{equation}
The composition equations
\[
  \WP{C_1;C_2}{Q}=\WP{C_1}{R},
  \qquad
  \WP{C_1;C_2}{g}=\WP{C_1}{h}
\]
rewrite the conclusion to \eqref{eq:cook-exact-bowtie}.

\paragraph{Case: $C$ is a conditional on $B$.}

The structurally smaller commands are $C_1$ and $C_2$.  Fix $g$ and define
\begin{equation}
  H\defsym[B]\cdot\WP{C_1}{g}
    +[\neg B]\cdot\WP{C_2}{g}.
  \label{eq:cook-if-H}
\end{equation}
By \eqref{eq:cook-wp-if},
\begin{equation}
  H=\WP{\ITE{B}{C_1}{C_2}}{g}.
  \label{eq:cook-if-H-wp}
\end{equation}
The indicator semantics gives
\begin{align}
  B&\models H=\WP{C_1}{g},
  \label{eq:cook-if-true}\\
  \neg B&\models H=\WP{C_2}{g}.
  \label{eq:cook-if-false}
\end{align}
These are semantic facts, not induction hypotheses or syntactic judgments.

For the arbitrary relation $\bowtie$, the two IHs are
\begin{align}
  &\vdash\HT{\top}{\WP{C_1}{g}}{C_1}{\top}{g},
  \label{eq:cook-if-IH-one}\\
  &\vdash\HT{\top}{\WP{C_2}{g}}{C_2}{\top}{g}.
  \label{eq:cook-if-IH-two}
\end{align}

\paragraph{The true branch by \textsc{Consequence}.}
Substitute
\[
  P=B,\quad f=H,\quad P'=\top,\quad f'=\WP{C_1}{g},
  \quad Q'=Q=\top,\quad g'=g
\]
in \eqref{eq:cook-consequence}.  The complete rule instance is
\begin{equation}
\frac{
  \begin{array}{c}
    B\models\top\land H\bowtie\WP{C_1}{g}
    \qquad
    \vdash\HT{\top}{\WP{C_1}{g}}{C_1}{\top}{g}
    \\[1mm]
    \top\models\top\land g\bowtie g
  \end{array}
}{
  \vdash\HT{B}{H}{C_1}{\top}{g}
}\;\textsc{Consequence}.
\label{eq:cook-if-conseq-true}
\end{equation}
The middle premise is IH~\eqref{eq:cook-if-IH-one}.  The first premise holds
because $B\models\top$ and \eqref{eq:cook-if-true} is an equality, hence it
implies each of $H\leq\WP{C_1}{g}$, $H=\WP{C_1}{g}$, and
$H\geq\WP{C_1}{g}$.  The last premise holds by reflexivity.  Thus the same
displayed derivation is valid for all three choices of $\bowtie$.

\paragraph{The false branch by \textsc{Consequence}.}
Now substitute
\[
  P=\neg B,\quad f=H,\quad P'=\top,
  \quad f'=\WP{C_2}{g},
  \quad Q'=Q=\top,\quad g'=g.
\]
The complete rule instance is
\begin{equation}
\frac{
  \begin{array}{c}
    \neg B\models\top\land H\bowtie\WP{C_2}{g}
    \qquad
    \vdash\HT{\top}{\WP{C_2}{g}}{C_2}{\top}{g}
    \\[1mm]
    \top\models\top\land g\bowtie g
  \end{array}
}{
  \vdash\HT{\neg B}{H}{C_2}{\top}{g}
}\;\textsc{Consequence}.
\label{eq:cook-if-conseq-false}
\end{equation}
The middle premise is IH~\eqref{eq:cook-if-IH-two}; the first follows from
the semantic equality \eqref{eq:cook-if-false}; and the last is reflexive.

\paragraph{Application of \textsc{If}.}
The conclusions of \eqref{eq:cook-if-conseq-true} and
\eqref{eq:cook-if-conseq-false} fill the two branch premises of
\textsc{If}:
\begin{equation}
\frac{
  \vdash\HT{B}{H}{C_1}{\top}{g}
  \qquad
  \vdash\HT{\neg B}{H}{C_2}{\top}{g}
}{
  \vdash\HT{\top}{H}{\ITE{B}{C_1}{C_2}}{\top}{g}
}\;\textsc{If}.
\label{eq:cook-if-derivation}
\end{equation}
The branch guards are $B\land\top=B$ and
$\neg B\land\top=\neg B$.  Finally,
\eqref{eq:cook-if-H-wp} rewrites $H$ to the exact pre-expectation.  This
proves the conditional case simultaneously for
$\bowtie\in\{\leq,=,\geq\}$.

\paragraph{Case: $C$ is a loop $L$ with guard $B$ and body $D$.}

\noindent\todo{TODO (equality case for while): Add an
equality-introduction rule combining the exact lower and upper loop judgments,
or provide a dedicated equality while rule. Until then, the derivations below
establish the two inequalities but do not yet construct the syntactic $[=]$
judgment.}\par

The only structurally smaller command is the body $D$.  Fix $g$ and define
\[
  W\defsym\WP{L}{g}.
\]
This is the semantic loop invariant: it is the exact expected value of $g$
after the whole loop.

\paragraph{Semantic fixed-point facts.}
Let
\[
  K_0\defsym0,
  \qquad
  K_{n+1}\defsym\Phi_g(K_n).
\]
By \eqref{eq:cook-loop-lfp} and Scott continuity,
\begin{equation}
  K_n\uparrow W,
  \qquad
  W=\Phi_g(W).
  \label{eq:cook-loop-kleene}
\end{equation}
Unfolding $\Phi_g$ and evaluating indicators gives the semantic equalities
\begin{align}
  B&\models W=\WP{D}{W},
  \label{eq:cook-loop-on-B}\\
  \neg B&\models W=g,
  \label{eq:cook-loop-on-not-B}\\
  B&\models K_{n+1}=\WP{D}{K_n}\qquad(n\geq0).
  \label{eq:cook-approximant-on-B}
\end{align}
None of these equalities is itself a syntactic Hoare derivation.

\paragraph{The body judgment, once for all three relations.}
The IH for $D$, instantiated at post-expectation $W$ and arbitrary relation
$\bowtie$, is
\begin{equation}
  \vdash\HT{\top}{\WP{D}{W}}{D}{\top}{W}.
  \label{eq:cook-loop-body-IH}
\end{equation}
Substitute
\[
  P=B,\quad f=W,\quad P'=\top,\quad f'=\WP{D}{W},
  \quad Q'=Q=\top,\quad g'=W
\]
in \textsc{Consequence}.  This gives the full rule instance
\begin{equation}
\frac{
  \begin{array}{c}
    B\models\top\land W\bowtie\WP{D}{W}
    \qquad
    \vdash\HT{\top}{\WP{D}{W}}{D}{\top}{W}
    \\[1mm]
    \top\models\top\land W\bowtie W
  \end{array}
}{
  \vdash\HT{B}{W}{D}{\top}{W}
}\;\textsc{Consequence}.
\label{eq:cook-loop-body-bowtie}
\end{equation}
The first premise follows from equality \eqref{eq:cook-loop-on-B}, the middle
premise is IH~\eqref{eq:cook-loop-body-IH}, and the last premise is reflexive.
The lower instance of \eqref{eq:cook-loop-body-bowtie} will be Premise~1 of
\textsc{While-lb}; its upper instance will be the body premise of
\textsc{While-ub}.

\paragraph{Residual witnesses for the lower loop rule.}
Because $K_n\uparrow W$, we have $0\leq K_n\leq W$.  Finiteness of $W$
allows us to define
\begin{equation}
  r_n\defsym W-K_n.
  \label{eq:cook-residual-definition}
\end{equation}
Each $r_n$ is a non-negative, finite-valued expectation.

\paragraph{Lower Premise 1.}
Instantiate \textsc{While-lb} with qualitative loop invariant $P:=\top$ and
quantitative loop invariant $W$.  Its first premise is
\[
  \vdash\HT[lower]{B}{W}{D}{\top}{W},
\]
which is exactly the $\bowtie={\leq}$ instance of
\eqref{eq:cook-loop-body-bowtie}. 

\paragraph{Lower Premise 2.}
The second premise is the semantic side condition
\[
  B\models r_0\geq W.
\]
Since $K_0=0$, definition \eqref{eq:cook-residual-definition} gives
$r_0=W-K_0=W$.  Hence $B\models r_0=W\geq W$.

\paragraph{Semantic residual equation for Lower Premise 3.}
For every $n$,
\begin{equation}
  B\models r_{n+1}=\WP{D}{r_n}.
  \label{eq:cook-residual-step}
\end{equation}
Indeed, $W=K_n+r_n$.  Additivity of $\mathsf{wp}$ gives
\[
  \WP{D}{W}=\WP{D}{K_n}+\WP{D}{r_n}.
\]
On $B$, equations \eqref{eq:cook-loop-on-B} and
\eqref{eq:cook-approximant-on-B} turn this into
\[
  W=K_{n+1}+\WP{D}{r_n}.
\]
Definition \eqref{eq:cook-residual-definition} also gives
$W=K_{n+1}+r_{n+1}$.  All terms are finite-valued, so cancellation yields
\eqref{eq:cook-residual-step}.  The equality appears because the chosen
$r_n=W-K_n$ is the \emph{exact} residual, not merely an arbitrary upper bound.
Premise~3 of \textsc{While-lb} only needs
$B\models r_{n+1}\geq\WP{D}{r_n}$, so
\eqref{eq:cook-residual-step} is a strictly stronger semantic fact.  It is not
yet a syntactic judgment; the required upper judgment is derived next.

\paragraph{Lower Premise 3.}
Fix an arbitrary $n$.  The upper IH for $D$, instantiated at post-expectation
$r_n$, is
\begin{equation}
  \vdash\HT[upper]{\top}{\WP{D}{r_n}}{D}{\top}{r_n}.
  \label{eq:cook-loop-residual-upper-IH}
\end{equation}
Substituting
\[
  P=B,\quad f=r_{n+1},\quad P'=\top,
  \quad f'=\WP{D}{r_n},
  \quad Q'=Q=\top,\quad g'=r_n,
  \quad\bowtie={\geq}
\]
in \textsc{Consequence} gives
\begin{equation}
\frac{
  \begin{array}{c}
    B\models\top\land r_{n+1}\geq\WP{D}{r_n}
    \qquad
    \vdash\HT[upper]{\top}{\WP{D}{r_n}}{D}{\top}{r_n}
    \\[1mm]
    \top\models\top\land r_n\geq r_n
  \end{array}
}{
  \vdash\HT[upper]{B}{r_{n+1}}{D}{\top}{r_n}
}\;\textsc{Consequence}.
\label{eq:cook-loop-residual-conseq}
\end{equation}
The first premise follows from \eqref{eq:cook-residual-step}, the middle one
is IH~\eqref{eq:cook-loop-residual-upper-IH}, and the last is reflexive.
This use of \textsc{Consequence} is essential: the IH has pre-guard $\top$ and
pre-expectation $\WP{D}{r_n}$, whereas the premise required by
\textsc{While-lb} has pre-guard $B$ and pre-expectation $r_{n+1}$.  The first
side condition performs exactly these two adaptations under the assumption
$B$.
Because $n$ was arbitrary, this derives
\[
  \forall n\geq0.\quad
  \vdash\HT[upper]{B}{r_{n+1}}{D}{\top}{r_n}.
\]

\paragraph{Lower Premise 4.}
The fourth premise is
\[
  B\models\lim_{n\to\infty}r_n=0.
\]
By \eqref{eq:cook-loop-kleene}, $K_n\uparrow W$.  Since $W$ is finite-valued,
pointwise
\[
  \lim_{n\to\infty}r_n
  =\lim_{n\to\infty}(W-K_n)
  =W-\lim_{n\to\infty}K_n
  =W-W=0.
\]
The equality holds on all states and thus on $B$.

\paragraph{Application of \textsc{While-lb}.}
The four established premises fill the lower loop rule as follows:
\begin{equation}
\frac{
  \begin{array}{c}
    \vdash\HT[lower]{B}{W}{D}{\top}{W}
    \qquad
    B\models r_0\geq W
    \\[1mm]
    \forall n.\ \vdash\HT[upper]{B}{r_{n+1}}{D}{\top}{r_n}
    \qquad
    B\models\lim_{n\to\infty}r_n=0
  \end{array}
}{
  \vdash\HT[lower]{\top}{W}{L}{\neg B}{W}
}\;\textsc{While-lb}.
\label{eq:cook-loop-lower-canonical}
\end{equation}

\paragraph{Application of \textsc{While-ub}.}
The $\bowtie={\geq}$ instance of \eqref{eq:cook-loop-body-bowtie} is
\[
  \vdash\HT[upper]{B}{W}{D}{\top}{W}.
\]
It is exactly the body premise obtained by taking $P=\top$ and $v=W$ in the
canonical upper loop rule:
\begin{equation}
\frac{
  \vdash\HT[upper]{B}{W}{D}{\top}{W}
}{
  \vdash\HT[upper]{\top}{W}{L}{\neg B}{W}
}\;\textsc{While-ub}.
\label{eq:cook-loop-upper-canonical}
\end{equation}

\paragraph{The common final consequence step.}
Equations \eqref{eq:cook-loop-lower-canonical} and
\eqref{eq:cook-loop-upper-canonical} establish, for the corresponding choice
of $\bowtie$,
\[
  \vdash\HT{\top}{W}{L}{\neg B}{W}.
\]
Both are converted to the desired post-assertion by the same fully
instantiated consequence rule:
\begin{equation}
\frac{
  \top\models\top\land W\bowtie W
  \qquad
  \vdash\HT{\top}{W}{L}{\neg B}{W}
  \qquad
  \neg B\models\top\land W\bowtie g
}{
  \vdash\HT{\top}{W}{L}{\top}{g}
}\;\textsc{Consequence}.
\label{eq:cook-loop-final-conseq}
\end{equation}
The first premise is reflexive.  The middle premise comes from
\textsc{While-lb} when $\bowtie={\leq}$ and from \textsc{While-ub} when
$\bowtie={\geq}$.  The last premise follows from
$\neg B\models W=g$ in \eqref{eq:cook-loop-on-not-B}; equality implies both
directions of $\bowtie$.  Since $W=\WP{L}{g}$ by definition, the conclusion
of \eqref{eq:cook-loop-final-conseq} is the exact-pre judgment for the loop.
This closes the lower and upper branches of the structural induction; the
equality branch is the explicit TODO recorded at the start of this case.
\end{proof}

\subsubsection{From the exact pre-expectation to every valid lower bound}

\begin{theorem}[Completeness for lower bounds]
\label{thm:cook-relative-completeness}
For every command $C$ and expectations $f,g$ in the expressive, finitary
fragment,
\[
  \valid\HT[lower]{\top}{f}{C}{\top}{g}
  \quad\Longrightarrow\quad
  \vdash\HT[lower]{\top}{f}{C}{\top}{g}.
\]
\end{theorem}

\begin{proof}
Assume
\[
  \valid\HT[lower]{\top}{f}{C}{\top}{g}.
\]
By semantic validity, the qualitative obligation
$\supp(C_m)\subseteq\top$ is automatic, while the quantitative obligation is
\[
  \forall m.\quad
  f(m)\leq\E{C,m}{[\top]\cdot g}.
\]
Because $[\top]=1$, definition \eqref{eq:cook-wp-definition} reduces it to
\begin{equation}
  f\leq\WP{C}{g}.
  \label{eq:cook-valid-bound}
\end{equation}
The $\bowtie={\leq}$ instance of the exact-pre lemma gives
\begin{equation}
  \vdash\HT[lower]{\top}{\WP{C}{g}}{C}{\top}{g}.
  \label{eq:cook-final-exact}
\end{equation}
Now substitute
\[
  P=P'=Q'=Q=\top,
  \quad f'=\WP{C}{g},
  \quad g'=g,
  \quad\bowtie={\leq}
\]
in \textsc{Consequence}.  The complete final rule instance is
\begin{equation}
\frac{
  \top\models\top\land f\leq\WP{C}{g}
  \qquad
  \vdash\HT[lower]{\top}{\WP{C}{g}}{C}{\top}{g}
  \qquad
  \top\models\top\land g\leq g
}{
  \vdash\HT[lower]{\top}{f}{C}{\top}{g}
}\;\textsc{Consequence}.
\label{eq:cook-final-consequence}
\end{equation}
The first premise is \eqref{eq:cook-valid-bound}, the middle premise is
\eqref{eq:cook-final-exact}, and the last premise is reflexive.  The
conclusion is exactly the desired derivability statement.
\end{proof}

\noindent\makebox[\linewidth]{\rule{\textwidth}{0.4pt}}
% ======================================================================
\section{Derived Rule}\label{app:derived-rule}
% ======================================================================

\paragraph{Multiplicative decay}\label{app:mult-decay}

\small\[
  \Infer[hoare][while-mult]
  {\HT[lower]{\pred}{\ex}{\WHILE\bexpr\DO\cmd}{\neg\bexpr\land\pred}{\ex}}
  {
    \HT[lower]{\bexpr\land\pred}{\ex}{\cmd}{\pred}{\ex}
    & \bexpr \land \pred \valid V \geq \ex
    & \HT[upper]{\bexpr\land\pred}{\beta \cdot V}{\cmd}{\top}{V}
    & \beta \in [0,1)
    & \bexpr \land \pred \valid V < \infty
  }
\]
\\
\subsection{Proof via Main Rule Soundness.}~\\
\textbf{Premises:}
\begin{align*}
  \text{(1) } [\bexpr \wedge \pred] \cdot \ex &\leq \WP{\cmd}{[\pred] \cdot \ex} \\
  \text{(2) } V &\geq \ex \\
  \text{(3) } \WP{\cmd}{V} &\leq \beta \cdot V, \quad \beta < 1 \\
  \text{(4) } V &< \infty
\end{align*}

\noindent\textbf{Proof Strategy:} $r_i \defsym \beta^i \cdot V$.

\noindent\textbf{Verification of Main Rule Premise 2:}
\begin{align*}
  r_0 = V \xrightarrow{\text{by (2)}} r_0 \ge \ex
    &\implies [\bexpr \wedge \pred] \cdot r_0 \ge [\bexpr \wedge \pred] \cdot \ex
\end{align*}

\noindent\textbf{Verification of Main Rule Premise 3:}
\begin{align*}
  \WP{\cmd}{r_i} &= \WP{\cmd}{\beta^i \cdot V} = \beta^i \cdot \WP{\cmd}{V} \\
    &\le \beta^i \cdot (\beta \cdot V) \quad \text{(by (3))} \\
    &= \beta^{i+1} \cdot V = r_{i+1} \\
  \implies [\bexpr \wedge \pred] \cdot r_{i+1} &\ge [\bexpr \wedge \pred] \cdot \WP{\cmd}{r_i}
\end{align*}

\noindent\textbf{Verification of Main Rule Premise 4:}
\begin{align*}
  \lim_{i \to \infty} r_i = \lim_{i \to \infty} (\beta^i \cdot V) \xrightarrow{\beta < 1, V < \infty} 0
    &\implies \lim_{i \to \infty} [\bexpr \wedge \pred] \cdot r_i = 0
\end{align*}

\noindent\textbf{Conclusion:} Soundness follows from the main rule.
\\
\subsection{Direct Proof.}~\\
\noindent\textbf{Definitions:} (Recall $I_n$ from \Cref{lem:In-recursion})
\begin{align*}
  r_i &\defsym \beta^i \cdot V
\end{align*}

\noindent\textbf{Auxiliary inequality on $\bexpr \wedge \pred$:}
\begin{align*}
  \WP{\cmd}{V} &\leq \beta \cdot V \quad \text{(premise 3)} \\
  \WP{\cmd}{r_n} = \beta^n \cdot \WP{\cmd}{V} &\leq \beta^{n+1} \cdot V = r_{n+1} \quad \text{(by linearity)} \\
  \WP{\cmd}{[\pred] \cdot r_n} \leq \WP{\cmd}{r_n} &\leq r_{n+1} \quad \text{(by monotonicity, } [\pred] \cdot r_n \leq r_n\text{)} \\
  [\bexpr \wedge \pred] \cdot r_{n+1} &\geq [\bexpr \wedge \pred] \cdot \WP{\cmd}{[\pred] \cdot r_n}
\end{align*}

\noindent\textbf{Key inequality:} $\forall n.\ [\pred] \cdot I_n + [\pred] \cdot r_n \geq [\pred] \cdot \ex$ (by induction)

\noindent\textit{Base case $n = 0$:}
\begin{align*}
  [\pred] \cdot I_0 + [\pred] \cdot r_0
    &= [\neg\bexpr \wedge \pred] \cdot \ex + [\pred] \cdot V \\
    &\geq [\neg\bexpr \wedge \pred] \cdot \ex + [\bexpr \wedge \pred] \cdot V \\
    &\geq [\neg\bexpr \wedge \pred] \cdot \ex + [\bexpr \wedge \pred] \cdot \ex \quad \text{(by premise 2)} \\
    &= [\pred] \cdot \ex
\end{align*}

\noindent\textit{Inductive step $n \to n+1$ (using \Cref{lem:In-recursion} for $I_{n+1}$):}
\begin{align*}
  [\pred] \cdot I_{n+1} + [\pred] \cdot r_{n+1}
    &= [\bexpr \wedge \pred] \cdot \WP{\cmd}{[\pred] \cdot I_n} + [\neg\bexpr \wedge \pred] \cdot \ex + [\pred] \cdot r_{n+1} \\
    &\geq [\bexpr \wedge \pred] \cdot \WP{\cmd}{[\pred] \cdot I_n} + [\neg\bexpr \wedge \pred] \cdot \ex + [\bexpr \wedge \pred] \cdot \WP{\cmd}{[\pred] \cdot r_n} \\
    &\quad\text{since } [\pred] \cdot r_{n+1} \geq [\bexpr \wedge \pred] \cdot r_{n+1} \geq [\bexpr \wedge \pred] \cdot \WP{\cmd}{[\pred] \cdot r_n} \\
    &\quad\text{(by } [\pred] \geq [\bexpr\wedge\pred] \text{ and auxiliary inequality)} \\
    &= [\bexpr \wedge \pred] \cdot \WP{\cmd}{[\pred] \cdot I_n + [\pred] \cdot r_n} + [\neg\bexpr \wedge \pred] \cdot \ex \quad \text{(by linearity)} \\
    &\geq [\bexpr \wedge \pred] \cdot \WP{\cmd}{[\pred] \cdot \ex} + [\neg\bexpr \wedge \pred] \cdot \ex \quad \text{(by IH + monotonicity)} \\
    &\geq [\bexpr \wedge \pred] \cdot \ex + [\neg\bexpr \wedge \pred] \cdot \ex \quad \text{(by premise 1)} \\
    &= [\pred] \cdot \ex
\end{align*}

\noindent\textbf{Limit argument:}
\begin{align*}
  \beta \in [0,1),\ \bexpr \wedge \pred \implies V < \infty
    &\implies \lim_{n \to \infty} r_n = 0 \\
  \lim_{n \to \infty} I_n
    &= \WP{\WHILE\bexpr\DO\cmd}{[\neg\bexpr \wedge \pred] \cdot \ex} \quad \text{(by Scott-continuity)} \\
  \implies [\pred] \cdot \ex &\leq \WP{\WHILE\bexpr\DO\cmd}{[\neg\bexpr \wedge \pred] \cdot \ex}
\end{align*}

% ----------------------------------------------------------------------
\noindent\todo{TODO (residual characterisation of \textsc{While-Mult}): The subsection characterising when \textsc{While-Mult} is instantiable
(geometric bound on the canonical residuals $H-I_n$) has been removed and
needs to be rewritten. Note that \texttt{prop:mult-residual} is still cited
in \texttt{benchmarks.tex}.}\par

% ----------------------------------------------------------------------
\subsection{Example: Geometric Termination}\label{app:mult-decay-geo}

\paragraph{Program.}
\[
  C_{\text{geo}} \;\defsym\; \WHILE b \DO \{ b <* \mathrm{Ber}(p); c <- c+1 \}
\]
At each iteration the loop resamples $b \sim \mathrm{Ber}(p)$ and increments the counter $c$.
The loop terminates when $\neg b$, which occurs geometrically with
$\mathbb{E}[T] = \tfrac{1}{1-p}$.

\paragraph{Instantiation.}
Set $B \equiv (b)$, $P \equiv \mathbf{true}$,
\[
  \ex(b,c) \;\defsym\; c + \frac{[b]}{1-p},
  \qquad
  V(b,c) \;\defsym\; \left(c + \frac{1}{1-p}\right)\cdot[b],
  \qquad
  \beta \;\defsym\; p(2-p).
\]
Note that $V = \ex$ on $B \wedge P$ (i.e.\ at $b$) and $V = 0$ otherwise.

\paragraph{Verification of the premises.}

\noindent\textbf{Premise~1} ($\HT[lower]{\bexpr\land\pred}{\ex}{\cmd}{\pred}{\ex}$):
\begin{align*}
  \WP{C_{\text{body}}}{\ex}(\mathbf{true},c)
    &= p\,\ex(\mathbf{true},c+1) + (1-p)\,\ex(\mathbf{false},c+1) \\
    &= p\!\left(c+1+\tfrac{1}{1-p}\right) + (1-p)(c+1)
     = c + \tfrac{1}{1-p} = \ex(\mathbf{true},c).
\end{align*}

\noindent\textbf{Premise~2} ($V \geq \ex$ on $B\land P$):
$V(\mathbf{true},c) = c + \tfrac{1}{1-p} = \ex(\mathbf{true},c)$.

\noindent\textbf{Premise~3} ($\WP{C_{\text{body}}}{V} \leq \beta\cdot V$ on $B\land P$):
\begin{align*}
  \WP{C_{\text{body}}}{V}(\mathbf{true},c)
    &= p\cdot V(\mathbf{true},c+1) + (1-p)\cdot V(\mathbf{false},c+1) \\
    &= p\!\left(c+1+\tfrac{1}{1-p}\right) = p\,V(\mathbf{true},c) + p.
\end{align*}
We need $p\,V + p \leq p(2-p)\,V$, i.e.\ $p \leq p(1-p)\,V = p(1-p)\!\left(c+\tfrac{1}{1-p}\right) = p\bigl((1-p)c+1\bigr)$.
Dividing by $p>0$: $1 \leq (1-p)c+1$, which reduces to $0 \leq (1-p)\,c$, true for all $c \geq 0$ (with equality at $c=0$).

\noindent\textbf{Premise~4} ($\beta \in [0,1)$):
$p(2-p) < 1 \iff (1-p)^2 > 0$, which holds for $p \in (0,1)$.

\noindent\textbf{Premise~5} ($V < \infty$ on $B\land P$):
$V(\mathbf{true},c) = c + \tfrac{1}{1-p} < \infty$ for any fixed $c$.

\paragraph{Witnesses and conclusion.}
The rule instantiates with $r_i(\mathbf{true},c) = \bigl(p(2-p)\bigr)^i\!\left(c+\tfrac{1}{1-p}\right)$,
giving
\[
  \vdash\; \HT[lower]{\mathbf{true}}{c + \tfrac{1}{1-p}}{C_{\text{geo}}}{\neg b}{c}.
\]
This certifies $\mathbb{E}[c_{\text{final}} \mid b=\mathbf{true},\, c=c_0] \geq c_0 + \tfrac{1}{1-p}$,
which matches the exact expected value of $c_0 + \mathbb{E}[T]$.


% \subsection{Additive decrease}\label{app:add-decrease}

% \paragraph{Rule.}
% \[
%   \Infer[hoare][while-add]
%   {\HT[lower]{\pred}{\ex}{\WHILE\bexpr\DO\cmd}{\neg\bexpr\land\pred}{\ex}}
%   {
%     \HT[lower]{\bexpr\land\pred}{\ex}{\cmd}{\pred}{\ex}
%     & \bexpr \land \pred \valid V \geq \ex
%     & [\bexpr \wedge \pred] \cdot V \geq \WP{\cmd}{[\pred] \cdot V} + \epsilon
%     & \epsilon > 0
%     & \bexpr \land \pred \valid V \leq M
%   }
% \]
% \\
% \subsubsection{Proof via Main Rule Soundness.}~\\
% \noindent\textbf{Premises:}
% \begin{align*}
%   \text{(1) } [\bexpr \wedge \pred] \cdot \ex &\le \WP{\cmd}{[\pred] \cdot \ex} \\
%   \text{(2) } V &\geq \ex \\
%   \text{(3) } \WP{\cmd}{[\pred] \cdot V} &\le V - \epsilon \\
%   \text{(4) } \epsilon > 0, \ &V \leq M
% \end{align*}

% \noindent\textbf{Proof Strategy:} $r_i \defsym \max(0, V - i \cdot \epsilon)$.

% \noindent\textbf{Verification of Main Rule Premise 2:}
% \begin{align*}
%   r_0 = V \xrightarrow{\text{by (2)}} r_0 \ge \ex
%     &\implies [\bexpr \wedge \pred] \cdot r_0 \ge [\bexpr \wedge \pred] \cdot \ex
% \end{align*}

% \noindent\textbf{Verification of Main Rule Premise 3:}
% \begin{align*}
%   \text{Assume } V - i\epsilon > 0 &\text{ (otherwise } r_i = 0 \text{ trivially):} \\
%   \WP{\cmd}{[\pred] \cdot r_i} &\le \WP{\cmd}{[\pred] \cdot (V - i\epsilon)} = \WP{\cmd}{[\pred] \cdot V} - i\epsilon \\
%     &\le (V - \epsilon) - i\epsilon \quad \text{(by (3))} \\
%     &= V - (i+1)\epsilon \le r_{i+1} \\
%   \implies [\bexpr \wedge \pred] \cdot r_{i+1} &\ge [\bexpr \wedge \pred] \cdot \WP{\cmd}{[\pred] \cdot r_i}
% \end{align*}

% \noindent\textbf{Verification of Main Rule Premise 4:}
% \begin{align*}
%   \text{Since } \epsilon > 0 \text{ and } V \leq M < \infty, &\ \exists N \text{ s.t. } \forall i > N: V - i\epsilon \le 0 \\
%   \therefore \lim_{i \to \infty} r_i &= 0 \implies \lim_{i \to \infty} [\bexpr \wedge \pred] \cdot r_i = 0
% \end{align*}

% \noindent\textbf{Conclusion:} Soundness follows from the main rule. \quad \checkmark
% \\
% \subsubsection{Direct Proof.}~\\

% \noindent\textbf{Definitions:} (Recall $I_n$ from \Cref{lem:In-recursion})
% \begin{align*}
%   \beta &\defsym 1 - \epsilon / M \\
%   r_i &\defsym \beta^i \cdot V
% \end{align*}

% \noindent\textbf{Range of $\beta$:}
% \begin{align*}
%   0 \leq \WP{\cmd}{[\pred] \cdot V} \leq V - \epsilon \leq M - \epsilon
%     &\implies \epsilon \leq M \\
%   0 < \epsilon \leq M
%     &\implies \beta \in [0, 1)
% \end{align*}

% \noindent\textbf{Auxiliary inequality on $\bexpr \wedge \pred$:}
% \begin{align*}
%   V \leq M
%     &\implies V \cdot \epsilon / M \leq \epsilon
%     \implies V - \epsilon \leq \beta \cdot V \\
%   \WP{\cmd}{[\pred] \cdot V} \leq V - \epsilon
%     &\implies \WP{\cmd}{[\pred] \cdot V} \leq \beta \cdot V \quad \text{(by premise 3)} \\
%   \beta^n \cdot \WP{\cmd}{[\pred] \cdot V}
%     &\leq \beta^{n+1} \cdot V \\
%   \WP{\cmd}{[\pred] \cdot r_n}
%     &\leq r_{n+1} \quad \text{(by linearity)} \\
%   \implies [\bexpr \wedge \pred] \cdot r_{n+1}
%     &\geq \WP{\cmd}{[\pred] \cdot r_n}
% \end{align*}

% \noindent\textbf{Key inequality:} $\forall n.\ [\pred] \cdot I_n + [\pred] \cdot r_n \geq [\pred] \cdot \ex$ (by induction)

% \noindent\textit{Base case $n = 0$:}
% \begin{align*}
%   [\pred] \cdot I_0 + [\pred] \cdot r_0
%     &= [\neg\bexpr \wedge \pred] \cdot \ex + [\pred] \cdot V \\
%     &\geq [\neg\bexpr \wedge \pred] \cdot \ex + [\bexpr \wedge \pred] \cdot V \\
%     &\geq [\neg\bexpr \wedge \pred] \cdot \ex + [\bexpr \wedge \pred] \cdot \ex \quad \text{(by premise 2)} \\
%     &= [\pred] \cdot \ex
% \end{align*}

% \noindent\textit{Inductive step $n \to n+1$ (using \Cref{lem:In-recursion} for $I_{n+1}$):}
% \begin{align*}
%   [\pred] \cdot I_{n+1} + [\pred] \cdot r_{n+1}
%     &= [\bexpr \wedge \pred] \cdot \WP{\cmd}{[\pred] \cdot I_n} + [\neg\bexpr \wedge \pred] \cdot \ex + [\pred] \cdot r_{n+1} \\
%     &\geq [\bexpr \wedge \pred] \cdot \WP{\cmd}{[\pred] \cdot I_n} + [\neg\bexpr \wedge \pred] \cdot \ex + [\bexpr \wedge \pred] \cdot \WP{\cmd}{[\pred] \cdot r_n} \\
%     &= [\bexpr \wedge \pred] \cdot \WP{\cmd}{[\pred] \cdot I_n + [\pred] \cdot r_n} + [\neg\bexpr \wedge \pred] \cdot \ex \quad \text{(by linearity)} \\
%     &\geq [\bexpr \wedge \pred] \cdot \WP{\cmd}{[\pred] \cdot \ex} + [\neg\bexpr \wedge \pred] \cdot \ex \quad \text{(by IH + monotonicity)} \\
%     &\geq [\bexpr \wedge \pred] \cdot \ex + [\neg\bexpr \wedge \pred] \cdot \ex \quad \text{(by premise 1)} \\
%     &= [\pred] \cdot \ex
% \end{align*}

% \noindent\textbf{Limit argument:}
% \begin{align*}
%   \beta \in [0,1),\ V \leq M
%     &\implies \lim_{n \to \infty} r_n = 0 \\
%   \lim_{n \to \infty} I_n
%     &= \WP{\WHILE\bexpr\DO\cmd}{[\neg\bexpr \wedge \pred] \cdot \ex} \quad \text{(by Scott-continuity)} \\
%   \implies [\pred] \cdot \ex &\leq \WP{\WHILE\bexpr\DO\cmd}{[\neg\bexpr \wedge \pred] \cdot \ex} \quad \checkmark
% \end{align*}

%%% Local Variables:
%%% mode: LaTeX
%%% TeX-master: "paper"
%%% End:



% ======================================================================
\section{Control Operators}%
\label{app:labeled-jumps}
% ======================================================================

\paragraph{Results.}
\[
  \ell\in\mathsf{Lab},
  \qquad
  \mathsf{Res}\ a::=\mathsf{Ok}(a)\mid\mathsf{Jump}(\ell,m),
\]
\[
  \mathsf{Ok}:a\to\mathsf{Res}\ a,
  \qquad
  \mathsf{Jump}:\mathsf{Lab}\to\Mem\to\mathsf{Res}\ a,
  \qquad
  \llbracket C\rrbracket:\Mem\to\mathcal D_{\leq1}(\mathsf{Res}\ \Mem).
\]
The missing probability mass represents non-termination.

\paragraph{Semantics.}
\begin{align*}
  \llbracket \vx <- \expr\rrbracket(m)
  &=\operatorname{return}
    \bigl(\mathsf{Ok}(m[\vx\mapped\expr(m)])\bigr),\\
  \llbracket\mathsf{jump}\ \ell\rrbracket(m)
  &=\operatorname{return}\bigl(\mathsf{Jump}(\ell,m)\bigr).
\end{align*}
\begin{equation}
  \llbracket C;D\rrbracket(m)
  =\operatorname{bind}\left(\llbracket C\rrbracket(m),\
  \lambda r.\,
  \begin{cases}
    \llbracket D\rrbracket(m'),&r=\mathsf{Ok}(m'),\\
    \operatorname{return}(\mathsf{Jump}(\ell,m')),
      &r=\mathsf{Jump}(\ell,m')
  \end{cases}\right)
  \label{eq:lj-sem-seq}
\end{equation}
\begin{equation}
  \llbracket\ell\{C\}\rrbracket(m)
  =\operatorname{bind}\left(\llbracket C\rrbracket(m),\
  \lambda r.\,
  \begin{cases}
    \operatorname{return}(\mathsf{Ok}(m')),
      &r=\mathsf{Ok}(m'),\\
    \operatorname{return}(\mathsf{Ok}(m')),
      &r=\mathsf{Jump}(\ell,m'),\\
    \operatorname{return}(\mathsf{Jump}(\ell',m')),
      &r=\mathsf{Jump}(\ell',m'),\ \ell'\neq\ell
  \end{cases}\right)
  \label{eq:lj-sem-label}
\end{equation}

\paragraph{Label environments.}
\[
  \mathsf{Pred}\defsym\Mem\to\Bool,
  \qquad
  \mathsf{Exp}\defsym\Mem\to[0,\infty].
\]
\[
  G:\mathsf{Lab}\to\mathsf{Pred},
  \qquad
  E:\mathsf{Lab}\to\mathsf{Exp}.
\]
The notation $G\mid E$ pairs the predicate and expectation
environments, as in a guarded assertion $\{P\mid f\}$.

\paragraph{Weakest pre-expectation and weakest guard.}
For a post-expectation $g:\mathsf{Exp}$ and a postcondition
$Q:\mathsf{Pred}$, the functions
$(E;\mathsf{Ok}\mapsto g):\mathsf{Res}\ \Mem\to[0,\infty]$ and
$(G;\mathsf{Ok}\mapsto Q):\mathsf{Res}\ \Mem\to\Bool$ read off a result with
$g$ (resp.\ $Q$) on a normal exit and with the environment on a jump:
\begin{equation}
  (E;\mathsf{Ok}\mapsto g)(r)
  \defsym
  \begin{cases}
    g(m'),&r=\mathsf{Ok}(m'),\\
    E(\ell)(m'),&r=\mathsf{Jump}(\ell,m'),
  \end{cases}
  \qquad
  (G;\mathsf{Ok}\mapsto Q)(r)
  \defsym
  \begin{cases}
    Q(m'),&r=\mathsf{Ok}(m'),\\
    G(\ell)(m'),&r=\mathsf{Jump}(\ell,m').
  \end{cases}
  \label{eq:ljwp-okmap}
\end{equation}
The two transformers $\mathrm{ewp}_{E}[C]:\mathsf{Exp}\to\mathsf{Exp}$
and $\mathrm{wp}_{G}[C]:\mathsf{Pred}\to\mathsf{Pred}$ are
\begin{equation}
  \mathrm{ewp}_{E}[C](g)(m)
  \defsym
  \mathbb E_{r\sim\llbracket C\rrbracket(m)}
  \bigl[(E;\mathsf{Ok}\mapsto g)(r)\bigr],
  \label{eq:ljwp-wp}
\end{equation}
\begin{equation}
  \mathrm{wp}_{G}[C](Q)(m)
  \defsym
  \forall r\in\supp(\llbracket C\rrbracket(m)).\
  (G;\mathsf{Ok}\mapsto Q)(r).
  \label{eq:ljwp-wpG}
\end{equation}
Nontermination is not in the support, so $\mathrm{wp}_{G}$ is liberal.
\paragraph{Semantic validity.}
\begin{equation}
  \begin{aligned}
  &\models G\mid E\vdash\{P\mid f\}_{\bowtie}\ C\ \{Q\mid g\}\\
  &\qquad\iff
  \forall m.\ P(m)\Longrightarrow
  \mathrm{wp}_{G}[C](Q)(m)
  \ \land\
  f(m)\mathrel{\bowtie}\mathrm{ewp}_{E}[C](g)(m).
  \end{aligned}
  \label{eq:ljwp-validity}
\end{equation}

\begin{lemma}[Basic equations]\label{lem:ljwp-equations}
For all $E,G,g,Q$, commands $C,D$ and labels $\ell$:
\begin{align}
  \mathrm{ewp}_{E}[\vx <- \expr](g)(m)&=g(m[\vx\mapped\expr(m)]),
  &\mathrm{wp}_{G}[\vx <- \expr](Q)(m)&=Q(m[\vx\mapped\expr(m)]),
  \label{eq:ljwp-assign}\\
  \mathrm{ewp}_{E}[\mathsf{jump}\ \ell](g)&=E(\ell),
  &\mathrm{wp}_{G}[\mathsf{jump}\ \ell](Q)&=G(\ell),
  \label{eq:ljwp-jump}\\
  \mathrm{ewp}_{E}[C;D](g)&=\mathrm{ewp}_{E}[C]\bigl(\mathrm{ewp}_{E}[D](g)\bigr),
  &\mathrm{wp}_{G}[C;D](Q)&=\mathrm{wp}_{G}[C]\bigl(\mathrm{wp}_{G}[D](Q)\bigr),
  \label{eq:ljwp-seq}\\
  \mathrm{ewp}_{E}[\ell\{C\}](g)&=\mathrm{ewp}_{E[\ell\mapsto g]}[C](g),
  &\mathrm{wp}_{G}[\ell\{C\}](Q)&=\mathrm{wp}_{G[\ell\mapsto Q]}[C](Q).
  \label{eq:ljwp-label}
\end{align}
\end{lemma}
\begin{proof}
Fix $m$. The proofs only use two facts about $\operatorname{return}$: for
every $v\in\mathsf{Res}\ \Mem$, $h:\mathsf{Res}\ \Mem\to[0,\infty]$ and
$\varphi:\mathsf{Res}\ \Mem\to\Bool$,
\begin{equation}
  \mathbb E_{r\sim\operatorname{return}(v)}[h(r)]=h(v),
  \qquad
  \bigl(\forall r\in\supp(\operatorname{return}(v)).\ \varphi(r)\bigr)
  \iff\varphi(v).
  \label{eq:lj-return}
\end{equation}
The first is \eqref{eq:cook-expect-return}, whose derivation does not
depend on $\Mem$ and therefore holds verbatim over $\mathsf{Res}\ \Mem$; the
second holds because $\supp(\operatorname{return}(v))=\{v\}$.

\smallskip\noindent\emph{Assignment.}
\[
  \begin{aligned}
  \mathrm{ewp}_{E}[\vx <- \expr](g)(m)
  &\overset{\eqref{eq:ljwp-wp}}{=}
  \mathbb E_{r\sim\llbracket \vx <- \expr\rrbracket(m)}
  \bigl[(E;\mathsf{Ok}\mapsto g)(r)\bigr]\\
  &\overset{\text{sem.}}{=}
  \mathbb E_{r\sim\operatorname{return}(\mathsf{Ok}(m[\vx\mapped\expr(m)]))}
  \bigl[(E;\mathsf{Ok}\mapsto g)(r)\bigr]\\
  &\overset{\eqref{eq:lj-return}}{=}
  (E;\mathsf{Ok}\mapsto g)\bigl(\mathsf{Ok}(m[\vx\mapped\expr(m)])\bigr)\\
  &\overset{\eqref{eq:ljwp-okmap}}{=}
  g(m[\vx\mapped\expr(m)]),
  \end{aligned}
\]
\[
  \begin{aligned}
  \mathrm{wp}_{G}[\vx <- \expr](Q)(m)
  &\overset{\eqref{eq:ljwp-wpG}}{\iff}
  \forall r\in\supp(\llbracket \vx <- \expr\rrbracket(m)).\
  (G;\mathsf{Ok}\mapsto Q)(r)\\
  &\overset{\text{sem.}}{\iff}
  \forall r\in\supp\bigl(\operatorname{return}(\mathsf{Ok}(m[\vx\mapped\expr(m)]))\bigr).\
  (G;\mathsf{Ok}\mapsto Q)(r)\\
  &\overset{\eqref{eq:lj-return}}{\iff}
  (G;\mathsf{Ok}\mapsto Q)\bigl(\mathsf{Ok}(m[\vx\mapped\expr(m)])\bigr)\\
  &\overset{\eqref{eq:ljwp-okmap}}{\iff}
  Q(m[\vx\mapped\expr(m)]).
  \end{aligned}
\]

\smallskip\noindent\emph{Jump.}
\[
  \begin{aligned}
  \mathrm{ewp}_{E}[\mathsf{jump}\ \ell](g)(m)
  &\overset{\eqref{eq:ljwp-wp}}{=}
  \mathbb E_{r\sim\llbracket \mathsf{jump}\ \ell\rrbracket(m)}
  \bigl[(E;\mathsf{Ok}\mapsto g)(r)\bigr]\\
  &\overset{\text{sem.}}{=}
  \mathbb E_{r\sim\operatorname{return}(\mathsf{Jump}(\ell,m))}
  \bigl[(E;\mathsf{Ok}\mapsto g)(r)\bigr]\\
  &\overset{\eqref{eq:lj-return}}{=}
  (E;\mathsf{Ok}\mapsto g)\bigl(\mathsf{Jump}(\ell,m)\bigr)\\
  &\overset{\eqref{eq:ljwp-okmap}}{=}
  E(\ell)(m),
  \end{aligned}
\]
\[
  \begin{aligned}
  \mathrm{wp}_{G}[\mathsf{jump}\ \ell](Q)(m)
  &\overset{\eqref{eq:ljwp-wpG}}{\iff}
  \forall r\in\supp(\llbracket \mathsf{jump}\ \ell\rrbracket(m)).\
  (G;\mathsf{Ok}\mapsto Q)(r)\\
  &\overset{\text{sem.}}{\iff}
  \forall r\in\supp\bigl(\operatorname{return}(\mathsf{Jump}(\ell,m))\bigr).\
  (G;\mathsf{Ok}\mapsto Q)(r)\\
  &\overset{\eqref{eq:lj-return}}{\iff}
  (G;\mathsf{Ok}\mapsto Q)\bigl(\mathsf{Jump}(\ell,m)\bigr)\\
  &\overset{\eqref{eq:ljwp-okmap}}{\iff}
  G(\ell)(m).
  \end{aligned}
\]
Neither $g$ nor $Q$ appears in the result, since $\mathsf{jump}\ \ell$
never terminates normally.

\smallskip\noindent\emph{Sequence and label.}
\todo{TODO: Proofs to be added.}
\end{proof}

\paragraph{Rules and their soundness.}

\paragraph{Assignment.}
For a predicate $Q$ and an expectation $g$ write
$Q[\vx\mapped\expr](m)\defsym Q(m[\vx\mapped\expr(m)])$ and
$g[\vx\mapped\expr](m)\defsym g(m[\vx\mapped\expr(m)])$.
\[
  \frac{}{
    G\mid E\vdash
    \{Q[\vx\mapped\expr]\mid g[\vx\mapped\expr]\}_{\bowtie}\
    \vx <- \expr\
    \{Q\mid g\}
  }
  \;\textsc{Assign}
\]

\begin{lemma}[\textsc{Assign}]\label{lem:ljwp-assign}
$\models G\mid E\vdash
\{Q[\vx\mapped\expr]\mid g[\vx\mapped\expr]\}_{\bowtie}\
\vx <- \expr\ \{Q\mid g\}$.
\end{lemma}
\noindent\todo{TODO: Proof to be rewritten in the new notation.}

\paragraph{Jump.}
A jump never ends normally, so the postcondition $\{Q\mid g\}$ plays no
role.
\[
  \frac{
    \forall m.\ P(m)\Longrightarrow G(\ell)(m)
    \qquad
    \forall m.\ P(m)\Longrightarrow
      f(m)\mathrel{\bowtie}E(\ell)(m)
  }{
    G\mid E\vdash
    \{P\mid f\}_{\bowtie}\ \mathsf{jump}\ \ell\ \{Q\mid g\}
  }
  \;\textsc{Jump}
\]

\begin{lemma}[\textsc{Jump}]\label{lem:ljwp-jump}
If both premises of \textsc{Jump} hold, then
$\models G\mid E\vdash\{P\mid f\}_{\bowtie}\ \mathsf{jump}\ \ell\ \{Q\mid g\}$.
\end{lemma}
\noindent\todo{TODO: Proof to be rewritten in the new notation.}

\paragraph{Sequence.}
The intermediate assertion $\{R\mid h\}$ connects the two premises: $R$
holds after every normal exit of $C$, and $h$ is the value carried from
$C$ to $D$ at that point. Jumps out of $C$ skip $D$ and are handled by the
shared environment $G\mid E$.
\[
  \frac{
    G\mid E\vdash\{P\mid f\}_{\bowtie}\ C\ \{R\mid h\}
    \qquad
    G\mid E\vdash\{R\mid h\}_{\bowtie}\ D\ \{Q\mid g\}
  }{
    G\mid E\vdash\{P\mid f\}_{\bowtie}\ C;D\ \{Q\mid g\}
  }
  \;\textsc{Seq}
\]

\begin{lemma}[\textsc{Seq}]\label{lem:ljwp-seq}
If $\models G\mid E\vdash\{P\mid f\}_{\bowtie}\ C\ \{R\mid h\}$ and
$\models G\mid E\vdash\{R\mid h\}_{\bowtie}\ D\ \{Q\mid g\}$, then
$\models G\mid E\vdash\{P\mid f\}_{\bowtie}\ C;D\ \{Q\mid g\}$.
\end{lemma}
\noindent\todo{TODO: Proof to be rewritten in the new notation.}

\paragraph{Label.}
Inside $\ell\{C\}$, a jump to $\ell$ ends the block normally. So in the
premise the label $\ell$ is given the postcondition itself: guard $Q$ and
weight $g$. Jumps to other labels keep their guard and weight from
$G\mid E$.
\[
  \frac{
    G[\ell\mapsto Q]\mid E[\ell\mapsto g]\vdash
    \{P\mid f\}_{\bowtie}\ C\ \{Q\mid g\}
  }{
    G\mid E\vdash
    \{P\mid f\}_{\bowtie}\ \ell\{C\}\ \{Q\mid g\}
  }
  \;\textsc{Label}
\]

\begin{lemma}[\textsc{Label}]\label{lem:ljwp-label}
If $\models G[\ell\mapsto Q]\mid E[\ell\mapsto g]\vdash
\{P\mid f\}_{\bowtie}\ C\ \{Q\mid g\}$, then
$\models G\mid E\vdash\{P\mid f\}_{\bowtie}\ \ell\{C\}\ \{Q\mid g\}$.
\end{lemma}
\noindent\todo{TODO: Proof to be rewritten in the new notation.}

\begin{theorem}[Soundness for labeled jumps]\label{thm:lj-soundness}
If $G\mid E\vdash\{P\mid f\}_{\bowtie}\ C\ \{Q\mid g\}$ is derivable with
the rules \textsc{Assign}, \textsc{Jump}, \textsc{Seq} and \textsc{Label},
then $\models G\mid E\vdash\{P\mid f\}_{\bowtie}\ C\ \{Q\mid g\}$.
\end{theorem}
\begin{proof}
By induction on the derivation. The last rule used is one of the four
rules above, and its lemma
(\Cref{lem:ljwp-assign,lem:ljwp-jump,lem:ljwp-seq,lem:ljwp-label}) shows
that the conclusion is valid, provided its premises are. For \textsc{Seq}
and \textsc{Label}, the premises are judgments, valid by the induction
hypothesis (for \textsc{Label}, in the environment
$G[\ell\mapsto Q]\mid E[\ell\mapsto g]$). For \textsc{Jump}, the premises
are conditions on memories that hold by assumption, and \textsc{Assign}
has no premises.
\end{proof}

\noindent\makebox[\linewidth]{\rule{\textwidth}{0.4pt}}

\section*{Canonical Witnesses: Worked Examples}\label{app:canonical-examples}

\paragraph{The canonical witness sequence.}
Given a loop $\WHILE\bexpr\DO\cmd$ with invariant $P$, postcondition expectation $g$, and
letting $q_n(m)=\Pr(T_m > n)$ denote the tail probabilities of the loop's
hitting time $T_m$ starting from state $m$, the \emph{canonical} witness sequence is
\begin{equation}\label{eq:canonical-r}
  r_n(m,c) \;\defsym\; (c + n + 1)\,q_n(m) \;+\; \sum_{j=n+1}^{\infty} q_j(m).
\end{equation}
This can equivalently be written as $r_n(m,c) = c\cdot q_n(m) + \mathbb{E}[T_m \cdot [T_m > n]]$,
which makes plain why it decreases to zero (the truncated expectation vanishes
as $n\to\infty$) and why it starts at $g$ (the Abel/tail-sum identity
$\sum_{n\geq 0} q_n(m) = \mathbb{E}[T_m]$ gives $r_0(m,c)=c+\mathbb{E}[T_m]=g(m,c)$).

We verify the four premises \textbf{P1}--\textbf{P4} of the \textsc{While-lb} rule on two concrete programs.

%----------------------------------------------------------------------
\subsection*{Example 1: Geometric Termination}

\paragraph{Program.}
\[
  C_{\text{geo}} \;\defsym\; \WHILE{b}{\;b \mathrel{:=} \mathrm{Ber}(p);\; c \mathrel{:=} c+1\;}
\]
The loop body samples a fresh Bernoulli($p$) bit and increments the counter $c$.
The loop terminates when $b=\mathbf{false}$, which happens geometrically:
$T_{b=\mathbf{true}} \sim \mathrm{Geom}(1-p)$, so $\mathbb{E}[T]=\tfrac{1}{1-p}$.

\paragraph{Postcondition and invariant.}
Take $P \equiv \mathbf{true}$ and
\[
  g(b,c) \;\defsym\; c + \frac{[b]}{1-p}.
\]
At the false-branch terminal state, $g(\mathbf{false},c)=c$, which records the
final value of $c$ (number of loop iterations taken).

\paragraph{Tail probabilities.}
$q_n(\mathbf{true}) = p^n$ and $q_n(\mathbf{false}) = 0$.

\paragraph{Canonical witnesses.}
\begin{align*}
  r_n(\mathbf{true},c) &= (c+n+1)\,p^n + \frac{p^{n+1}}{1-p}, \\
  r_n(\mathbf{false},c) &= 0.
\end{align*}

\paragraph{Verification of P1.}
We need $\vdash_{\leq}\HT{b\land P}{g}{C_{\text{body}}}{P}{g}$, i.e.\
$\mathrm{wp}[b\mathrel{:=}\mathrm{Ber}(p);\,c\mathrel{:=}c+1](g) \geq g$.
\begin{align*}
  \mathrm{wp}[C_{\text{body}}](g)(b,c)
    &= p\cdot g(\mathbf{true},c+1) + (1-p)\cdot g(\mathbf{false},c+1) \\
    &= p\!\left(c+1+\tfrac{1}{1-p}\right) + (1-p)(c+1) \\
    &= c + 1 + \tfrac{p}{1-p}
     = c + \tfrac{1}{1-p}
     = g(\mathbf{true},c).
\end{align*}
Equality holds, so P1 is satisfied.

\paragraph{Verification of P2.}
We need $r_0 \geq g$ on $B\land P$, i.e.\ at $b=\mathbf{true}$:
\[
  r_0(\mathbf{true},c)
  = (c+0+1)\,p^0 + \frac{p^1}{1-p}
  = c + 1 + \frac{p}{1-p}
  = c + \frac{1}{1-p}
  = g(\mathbf{true},c).
\]

\paragraph{Verification of P3.}
We need $\mathrm{wp}[C_{\text{body}}](r_n) \leq r_{n+1}$. Compute at $b=\mathbf{true}$:
\begin{align*}
  \mathrm{wp}[C_{\text{body}}](r_n)(\mathbf{true},c)
    &= p\cdot r_n(\mathbf{true},c+1) + (1-p)\cdot r_n(\mathbf{false},c+1) \\
    &= p\!\left[(c+n+2)\,p^n + \frac{p^{n+1}}{1-p}\right] + 0 \\
    &= (c+n+2)\,p^{n+1} + \frac{p^{n+2}}{1-p}.
\end{align*}
Meanwhile,
\[
  r_{n+1}(\mathbf{true},c)
  = (c+n+2)\,p^{n+1} + \frac{p^{n+2}}{1-p}.
\]
These are equal, so P3 holds with equality.

\paragraph{Verification of P4.}
On $B\land P$ (i.e.\ $b=\mathbf{true}$):
\[
  \lim_{n\to\infty} r_n(\mathbf{true},c)
  = \lim_{n\to\infty}\!\left[(c+n+1)\,p^n + \frac{p^{n+1}}{1-p}\right] = 0,
\]
since $p\in(0,1)$ so exponential decay dominates the polynomial prefactor.

%----------------------------------------------------------------------
\subsection*{Example 2: Deterministic Countdown}

\paragraph{Program.}
\[
  C_{\text{det}} \;\defsym\; \WHILE{x>0}{\;x \mathrel{:=} x-1;\; c \mathrel{:=} c+1\;}
\]
The loop body decrements $x$ and increments $c$ deterministically.
Starting from $x=k\geq 0$, the loop runs exactly $k$ iterations: $T_k = k$.

\paragraph{Postcondition and invariant.}
Take $P \equiv \mathbf{true}$ and
\[
  g(x,c) \;\defsym\; c + x.
\]

\paragraph{Tail probabilities.}
Because $T_x = x$ deterministically,
$q_n(x) = \Pr(T_x > n) = [x > n]$.

\paragraph{Canonical witnesses.}
\begin{equation*}
  r_n(x,c) = (c+n+1)\,[x>n] + \sum_{j=n+1}^{\infty}[x>j]
            = (c+n+1)\,[x>n] + \max(x-n-1,\,0).
\end{equation*}
For $x > n$ this simplifies to $(c+n+1) + (x-n-1) = c+x$.
For $x \leq n$ both terms vanish, giving $r_n(x,c) = 0$.
Hence
\[
  r_n(x,c) = (c + x)\cdot[x > n].
\]

\paragraph{Verification of P1.}
\begin{align*}
  \mathrm{wp}[C_{\text{body}}](g)(x,c)
    &= g(x-1,\,c+1) = (c+1)+(x-1) = c+x = g(x,c).
\end{align*}

\paragraph{Verification of P2.}
At $x>0$ (i.e.\ $B\land P$):
\[
  r_0(x,c) = (c+x)\cdot[x>0] = c+x = g(x,c).
\]

\paragraph{Verification of P3.}
\begin{align*}
  \mathrm{wp}[C_{\text{body}}](r_n)(x,c)
    &= r_n(x-1,\,c+1) \\
    &= (c+1+(x-1))\cdot[x-1>n] \\
    &= (c+x)\cdot[x>n+1] \\
    &= r_{n+1}(x,c).
\end{align*}

\paragraph{Verification of P4.}
For any fixed $x$, $[x>n]=0$ for all $n\geq x$, so
$\lim_{n\to\infty} r_n(x,c) = 0$.

% % ======================================================================
% Benchmark examples for the While-lb rule and the derived While-Mult rule.
% This file is self-contained and \input at the end of paper.tex.
% ======================================================================
\newpage

% The style file makes `<' active (it renders as \le, or as an assignment
% arrow when followed by - or *).  \lt gives a genuine strict inequality.
{\catcode`\<=12 \gdef\lt{\mathrel{<}}}

% ======================================================================
\section{Benchmark Examples}\label{app:benchmarks}
% ======================================================================

We validate the general rule \textsc{While-lb} of \Cref{fig:hoare:while-lb}
and the derived rule \textsc{While-Mult} of \Cref{app:derived-rule} on five
benchmarks from the literature on lower bounds for probabilistic loops.
Throughout, $\cmd \oplus_{p} \cmdtwo$ denotes probabilistic choice,
\[
  \WP{\cmd \oplus_{p} \cmdtwo}{f}
  \;=\;
  p \cdot \WP{\cmd}{f} + (1-p)\cdot\WP{\cmdtwo}{f},
\]
where the probability $p$ may be a state-dependent expression. For a loop
$\WHILE\bexpr\DO\cmd$ and a memory $\mem$, $T_\mem$ denotes its termination
time (number of body executions) started from $\mem$, and
$q_n(\mem) \defsym \Pr(T_\mem > n)$ its tail probabilities, as in
\Cref{app:canonical-examples}. In every example below the certified lower
bound coincides with the exact quantity, i.e.\ the rules are applied at full
tightness.

\paragraph{Choosing between the two rules.}
The witnesses of \textsc{While-Mult} are geometric, $r_i = \beta^i V$, and this
is not an accident of the proof: whenever the rule applies, the mass still
trapped in the loop after $n$ iterations \emph{must} decay geometrically. The
following proposition makes the criterion precise; we use it below to show
that two of the five benchmarks provably lie outside the scope of
\textsc{While-Mult}.

\begin{proposition}[Geometric tails are necessary for \textsc{While-Mult}]%
\label{prop:geo-tail}
Suppose premises 1--5 of \textsc{While-Mult} hold for the termination triple
$\HT[lower]{\pred}{1}{\WHILE\bexpr\DO\cmd}{\neg\bexpr\land\pred}{1}$.
Then for every $\mem \vDash \bexpr\land\pred$ and every $n \geq 0$,
\[
  \Pr(T_\mem > n) \;\leq\; \beta^{n}\, V(\mem).
\]
\end{proposition}

\begin{proof}
As in \Cref{app:derived-rule}, the sequence $r_n \defsym \beta^n V$ satisfies
premises 2--4 of the main rule, so \Cref{lem:key-ineq} yields
$[\pred]\cdot I_n + [\pred]\cdot r_n \geq [\pred]\cdot 1$ with
$I_n = [\pred]\cdot\WP{W_n}{[\pred\wedge\neg\bexpr]\cdot 1}$. Since $W_n$
aborts all runs exceeding $n$ iterations,
$I_n(\mem) \leq \Pr(T_\mem \leq n)$, and therefore, at any
$\mem \vDash \bexpr \land \pred$,
\[
  \Pr(T_\mem > n)
  \;=\; 1 - \Pr(T_\mem \leq n)
  \;\leq\; 1 - I_n(\mem)
  \;\leq\; r_n(\mem)
  \;=\; \beta^n V(\mem).
  \qedhere
\]
\end{proof}

Consequently, a loop that terminates almost surely but whose termination time
has a polynomial tail admits \emph{no} instantiation of \textsc{While-Mult}
for the termination triple, no matter how $V$ and $\beta$ are chosen
(\Cref{app:bench-spline,app:bench-rw}). Conversely, geometric tails do not by
themselves make the instantiation obvious: the natural candidate $V = \ex$
may fail and a reweighted $V$ is needed
(\Cref{app:bench-nb,app:bench-cc}).

\begin{table}[H]
  \centering
  \small
  \begin{tabular}{lllll}
    \toprule
    Benchmark & Certified bound & Tail of $T$ & Rule \\
    \midrule
    Gambler's ruin (\ref{app:bench-ruin})
      & $\Pr(\text{win}) \geq x/N$
      & $\Theta(\cos^n(\pi/N))$
      & \textsc{While-Mult} \\
    Negative binomial (\ref{app:bench-nb})
      & $\mathbb{E}[c_{\text{final}}] \geq c + x/p$
      & geometric
      & \textsc{While-Mult} \\
    Coupon collector (\ref{app:bench-cc})
      & $\mathbb{E}[c_{\text{final}}] \geq c + N\mathcal{H}_k$
      & $\Theta((1-\sfrac{1}{N})^n)$
      & \textsc{While-Mult} \\
    Escaping spline (\ref{app:bench-spline})
      & $\Pr(\text{termination}) \geq 1$
      & $x/(x+n)$
      & \textsc{While-lb} \\
    Symmetric random walk (\ref{app:bench-rw})
      & $\Pr(\text{termination}) \geq 1$
      & $\Theta(x/\sqrt{n})$
      & \textsc{While-lb} \\
    \bottomrule
  \end{tabular}
  \caption{The five benchmarks. For the last two, \textsc{While-Mult} is
  provably inapplicable by \Cref{prop:geo-tail}.}
\end{table}

% ----------------------------------------------------------------------
\subsection{Gambler's Ruin (Bounded Symmetric Random Walk)}%
\label{app:bench-ruin}
% ----------------------------------------------------------------------

\paragraph{Program.}
Fix an integer $N \geq 2$ and let $x$ range over $\{0,\dots,N\}$:
\[
  C_{\text{ruin}} \;\defsym\;
  \WHILE{0 \lt x \lt N}{\;
    (x \mathrel{:=} x-1) \oplus_{1/2} (x \mathrel{:=} x+1)
  \;}
\]
A fair gambler starts with $x$ coins and plays until ruin ($x = 0$) or
victory ($x = N$). The loop terminates almost surely and the classical
winning probability is exactly $x/N$.

\paragraph{Specification.}
Take $\pred \equiv 0 \leq x \leq N$, $\bexpr \equiv 0 \lt x \lt N$, and
\[
  \ex(x) \;\defsym\; \frac{x}{N}.
\]
On terminal states $x \in \{0, N\}$ the expectation $\ex$ equals the winning
indicator $[x = N]$, so the target triple
$\HT[lower]{\pred}{\ex}{C_{\text{ruin}}}{\neg\bexpr\land\pred}{\ex}$
certifies $\Pr(x_{\text{final}} = N) \geq x_0/N$.

\paragraph{Instantiation of \textsc{While-Mult}.}
\[
  V(x) \;\defsym\; N \sin\!\Bigl(\frac{\pi x}{N}\Bigr)
  \quad (0 \leq x \leq N; \text{ and } V \defsym 0 \text{ elsewhere}),
  \qquad
  \beta \;\defsym\; \cos\!\Bigl(\frac{\pi}{N}\Bigr).
\]

\paragraph{Premise 1} ($\HT[lower]{\bexpr\land\pred}{\ex}{\cmd}{\pred}{\ex}$):
both successor states lie in $[0,N]$, so the $[\pred]$ factor is $1$ and
\[
  \WP{C_{\text{body}}}{\ex}(x)
  = \frac{1}{2}\cdot\frac{x-1}{N} + \frac{1}{2}\cdot\frac{x+1}{N}
  = \frac{x}{N}
  = \ex(x).
  \checkmark
\]
Equality holds: $\ex$ is harmonic for the fair walk.

\paragraph{Premise 2} ($V \geq \ex$ on $\bexpr\land\pred$):
by Jordan's inequality $\sin\theta \geq \tfrac{2}{\pi}\theta$ on
$[0,\sfrac{\pi}{2}]$ and the symmetry
$\sin(\pi x/N) = \sin(\pi(N-x)/N)$,
\[
  V(x) \;\geq\; 2\min(x,\,N-x) \;\geq\; 2 \;>\; 1 \;\geq\; \ex(x)
  \qquad (1 \leq x \leq N-1).
  \checkmark
\]

\paragraph{Premise 3} ($\WP{C_{\text{body}}}{V} \leq \beta V$ on
$\bexpr\land\pred$): by the sum-to-product identity
$\sin(a-d) + \sin(a+d) = 2\sin a \cos d$,
\[
  \WP{C_{\text{body}}}{V}(x)
  = \frac{N}{2}\Bigl(\sin\tfrac{\pi(x-1)}{N} + \sin\tfrac{\pi(x+1)}{N}\Bigr)
  = N\sin\tfrac{\pi x}{N}\,\cos\tfrac{\pi}{N}
  = \beta\, V(x).
  \checkmark
\]
Equality holds even at the boundary cases $x = 1$ and $x = N-1$, since
$V(0) = V(N) = 0$: $V$ is an eigenfunction of the one-step operator of the
walk killed at $\{0, N\}$.

\paragraph{Premises 4 and 5.}
$\beta = \cos(\pi/N) \in [0,1)$ for $N \geq 2$, and $V \leq N \lt \infty$.
$\checkmark$

\paragraph{Conclusion.}
With witnesses $r_i(x) = \cos^i(\pi/N)\cdot N\sin(\pi x/N)$,
\[
  \vdash\; \HT[lower]{0 \leq x \leq N}{\tfrac{x}{N}}
     {C_{\text{ruin}}}{x \in \{0,N\}}{\tfrac{x}{N}},
\]
i.e.\ $\Pr(\text{the gambler wins}) \geq x_0/N$, matching the exact value.
Moreover $\beta$ is optimal: $\cos(\pi/N)$ is the spectral radius of the
killed walk, and indeed $q_n(x) = \Theta(\cos^n(\pi/N))$, so by
\Cref{prop:geo-tail} no instantiation can certify a faster decay.

\paragraph{Variant: biased walk.}
For $(x \mathrel{:=} x+1) \oplus_{p} (x \mathrel{:=} x-1)$ with
$p \in (0,1)$, $q \defsym 1-p$ and $\lambda \defsym q/p \neq 1$, the same
scheme applies with
\[
  \ex(x) = \frac{\lambda^x - 1}{\lambda^N - 1},
  \qquad
  V(x) = A\,\lambda^{x/2} \sin\!\Bigl(\frac{\pi x}{N}\Bigr),
  \qquad
  \beta = 2\sqrt{pq}\,\cos\!\Bigl(\frac{\pi}{N}\Bigr),
\]
where $A \defsym \tfrac{N}{2}\max\bigl(1, \lambda^{-(N-1)/2}\bigr)$.
Premise~1 holds with equality since
$p\lambda^{x+1} + q\lambda^{x-1} = \lambda^{x}(q + p) = \lambda^x$;
Premise~3 holds with equality since
$p\,\lambda^{1/2} = q\,\lambda^{-1/2} = \sqrt{pq}$, whence
$p V(x+1) + q V(x-1)
 = \lambda^{x/2}A\sqrt{pq}\,\bigl(\sin\tfrac{\pi(x+1)}{N}
 + \sin\tfrac{\pi(x-1)}{N}\bigr) = \beta V(x)$;
Premise~2 follows as before from
$A\lambda^{x/2} \geq \tfrac{N}{2}$ on $1 \leq x \leq N-1$, giving
$V(x) \geq 1 \geq \ex(x)$; and
$\beta \lt 1$ since $2\sqrt{pq} \leq 1$ by AM--GM (strict for
$p \neq \tfrac12$) and $\cos(\pi/N) \lt 1$.

% ----------------------------------------------------------------------
\subsection{Negative Binomial}\label{app:bench-nb}
% ----------------------------------------------------------------------

\paragraph{Program.}
Fix $p \in (0,1]$ and let $x, c$ range over the naturals:
\[
  C_{\text{nb}} \;\defsym\;
  \WHILE{x > 0}{\;
    \bigl((x \mathrel{:=} x-1) \oplus_{p} \SKIP\bigr);\;
    c \mathrel{:=} c+1
  \;}
\]
Each iteration succeeds with probability $p$; the loop performs $x$
successes, so $T$ is negative binomial with $\mathbb{E}[T] = x/p$.

\paragraph{Specification.}
Take $\pred \equiv x \geq 0 \land c \geq 0$, $\bexpr \equiv x > 0$, and
\[
  \ex(x,c) \;\defsym\; c + \frac{x}{p}.
\]
The target triple certifies
$\mathbb{E}[c_{\text{final}}] \geq c_0 + x_0/p$.

\paragraph{Why the naive $V$ fails.}
The candidate $V = \ex$ does \emph{not} satisfy Premise~3: for $x \geq 2$ no
mass leaves the loop within one iteration and (see Premise~1 below)
$\WP{C_{\text{body}}}{\ex} = \ex$, so $\WP{C_{\text{body}}}{\ex} \leq
\beta\,\ex$ with $\beta \lt 1$ would force $\ex \leq 0$ --- impossible, as
$\ex > 0$ on $\bexpr$. This is not a defect of the loop (its tail \emph{is}
geometric) but of the candidate: the remedy is to reweight $V$ exponentially
in $x$, trading a larger dominating function for a uniform one-step
contraction.

\paragraph{Instantiation of \textsc{While-Mult}.}
Abbreviating $b_x \defsym (x+2)/p$,
\[
  V(x,c) \;\defsym\; [x \geq 1]\cdot 2^{x}\,(c + b_x),
  \qquad
  \beta \;\defsym\; 1 - \frac{p}{4}.
\]

\paragraph{Premise 1} (with equality; successors satisfy $\pred$):
\[
  \WP{C_{\text{body}}}{\ex}(x,c)
  = p\Bigl(c+1+\frac{x-1}{p}\Bigr) + (1-p)\Bigl(c+1+\frac{x}{p}\Bigr)
  = c + 1 + \frac{x}{p} - 1
  = \ex(x,c).
  \checkmark
\]

\paragraph{Premise 2} ($V \geq \ex$ on $\bexpr\land\pred$):
$2^x \geq 1$ and $b_x \geq x/p$. $\checkmark$

\paragraph{Premise 3} ($\WP{C_{\text{body}}}{V} \leq \beta V$ on
$\bexpr\land\pred$): we distinguish whether the success branch stays in the
loop.

\emph{Case $x \geq 2$.} Using $b_{x-1} = b_x - 1/p$,
\begin{align*}
  \WP{C_{\text{body}}}{V}(x,c)
  &= p\,2^{x-1}(c+1+b_{x-1}) + (1-p)\,2^{x}(c+1+b_x) \\
  &= 2^{x}\Bigl[\tfrac{p}{2}\bigl(c+1+b_x - \tfrac1p\bigr)
       + (1-p)(c+1+b_x)\Bigr]
   = 2^{x}\Bigl[\bigl(1-\tfrac{p}{2}\bigr)(c+1+b_x) - \tfrac12\Bigr].
\end{align*}
The required bound $\leq \bigl(1-\tfrac{p}{4}\bigr)2^x(c+b_x)$ rearranges to
\[
  \frac{1-p}{2} \;\leq\; \frac{p}{4}\,(c + b_x),
\]
which holds since $c \geq 0$ and $\tfrac{p}{4}\, b_x \geq \tfrac{p}{4}\cdot
\tfrac{4}{p} = 1 > \tfrac{1-p}{2}$. $\checkmark$

\emph{Case $x = 1$} (success exits the loop; $V(0,\cdot) = 0$, $b_1 = 3/p$):
\[
  \WP{C_{\text{body}}}{V}(1,c)
  = (1-p)\cdot 2\,(c+1+\tfrac{3}{p})
  \;\overset{!}{\leq}\;
  \bigl(1-\tfrac{p}{4}\bigr)\cdot 2\,(c+\tfrac{3}{p}).
\]
Comparing coefficients of $c$: $1-p \leq 1-\tfrac{p}{4}$. Comparing
constants:
$(1-p)\bigl(1+\tfrac{3}{p}\bigr) \leq \tfrac{3}{p} - \tfrac{3}{4}
 \iff 1-p \leq \tfrac{9}{4}$,
which always holds. $\checkmark$

\paragraph{Premises 4 and 5.}
$\beta = 1 - p/4 \in [\sfrac34, 1)$ for $p \in (0,1]$, and $V$ is pointwise
finite. $\checkmark$

\paragraph{Conclusion.}
\[
  \vdash\; \HT[lower]{x \geq 0 \land c \geq 0}{c + \tfrac{x}{p}}
     {C_{\text{nb}}}{x = 0 \land c \geq 0}{c + \tfrac{x}{p}},
\]
certifying $\mathbb{E}[c_{\text{final}}] \geq c_0 + x_0/p$ --- the exact
expected value. The certified decay rate $\beta = 1-p/4$ is coarser than the
true tail rate $1-p$; this is harmless, as \textsc{While-Mult} only requires
\emph{some} valid pair $(V,\beta)$, not a sharp one.

% ----------------------------------------------------------------------
\subsection{Coupon Collector}\label{app:bench-cc}
% ----------------------------------------------------------------------

\paragraph{Program.}
We track $k \in \{0,\dots,N\}$, the number of coupons still missing, and the
draw counter $c$; a uniform draw hits a missing coupon with probability
$k/N$:
\[
  C_{\text{cc}} \;\defsym\;
  \WHILE{k > 0}{\;
    \bigl((k \mathrel{:=} k-1) \oplus_{k/N} \SKIP\bigr);\;
    c \mathrel{:=} c+1
  \;}
\]
Writing $\mathcal{H}_k \defsym \sum_{j=1}^{k} \sfrac{1}{j}$, the expected
number of draws from $k$ missing coupons is $N \mathcal{H}_k$; from a fresh
start ($k = N$) it is the classical $N \mathcal{H}_N$.

\paragraph{Specification.}
Take $\pred \equiv 0 \leq k \leq N \land c \geq 0$, $\bexpr \equiv k > 0$,
and
\[
  \ex(k,c) \;\defsym\; c + N \mathcal{H}_k.
\]

\paragraph{Tail behaviour.}
The tail is genuinely geometric with rate $1 - \sfrac1N$: fixing any single
missing coupon gives $\Pr(T > n) \geq (1-\sfrac1N)^n$, and a union bound over
the $k$ missing coupons gives $\Pr(T > n) \leq k\,(1-\sfrac1N)^n$. So
\textsc{While-Mult} is the right tool; as in \Cref{app:bench-nb} the naive
$V = \ex$ fails (for $k \geq 2$ no mass exits within one step), and we
reweight exponentially in $k$.

\paragraph{Instantiation of \textsc{While-Mult}.}
Abbreviating $b_k \defsym N\mathcal{H}_k + N$ (so $b_{k-1} = b_k - N/k$),
\[
  V(k,c) \;\defsym\; [k \geq 1]\cdot 2^{k}\,(c + b_k),
  \qquad
  \beta \;\defsym\; 1 - \frac{1}{2N}.
\]

\paragraph{Premise 1} (with equality; successors satisfy $\pred$):
\[
  \WP{C_{\text{body}}}{\ex}(k,c)
  = \frac{k}{N}\bigl(c+1+N\mathcal{H}_{k-1}\bigr)
  + \Bigl(1-\frac{k}{N}\Bigr)\bigl(c+1+N\mathcal{H}_k\bigr)
  = c + 1 + N\mathcal{H}_k - \frac{k}{N}\cdot\frac{N}{k}
  = \ex(k,c).
  \checkmark
\]

\paragraph{Premise 2.} $V \geq \ex$ on $k \geq 1$: immediate. $\checkmark$

\paragraph{Premise 3} ($\WP{C_{\text{body}}}{V} \leq \beta V$ on
$\bexpr\land\pred$):

\emph{Case $k \geq 2$} (both successors in the loop):
\begin{align*}
  \WP{C_{\text{body}}}{V}(k,c)
  &= \frac{k}{N}\, 2^{k-1}\Bigl(c+1+b_k-\frac{N}{k}\Bigr)
   + \Bigl(1-\frac{k}{N}\Bigr) 2^{k}\,(c+1+b_k) \\
  &= 2^{k}\Bigl[\Bigl(1-\frac{k}{2N}\Bigr)(c+1+b_k) - \frac12\Bigr].
\end{align*}
The required bound $\leq \bigl(1-\tfrac{1}{2N}\bigr)2^k(c+b_k)$ rearranges to
\[
  \frac12 - \frac{k}{2N}
  \;\leq\;
  \frac{k-1}{2N}\,(c+b_k),
\]
which holds since $k \geq 2$, $c \geq 0$ and
$\tfrac{k-1}{2N}\,b_k \geq \tfrac{1}{2N}\bigl(N\mathcal{H}_2+N\bigr)
 = \tfrac54 \geq \tfrac12$. $\checkmark$

\emph{Case $k = 1$} (success exits the loop; $V(0,\cdot) = 0$, $b_1 = 2N$):
\[
  \WP{C_{\text{body}}}{V}(1,c)
  = \Bigl(1-\frac1N\Bigr)\cdot 2\,(c+1+2N)
  \;\overset{!}{\leq}\;
  \Bigl(1-\frac{1}{2N}\Bigr)\cdot 2\,(c+2N).
\]
Coefficients of $c$: $1-\sfrac1N \leq 1-\sfrac{1}{2N}$. Constants:
$(1-\sfrac1N)(2N+1) = 2N-1-\sfrac1N \leq 2N-1
 = (1-\sfrac{1}{2N})\cdot 2N$. $\checkmark$

\paragraph{Premises 4 and 5.}
$\beta = 1-\sfrac{1}{2N} \in [0,1)$ and $V$ is pointwise finite.
$\checkmark$

\paragraph{Conclusion.}
\[
  \vdash\; \HT[lower]{\pred}{c + N\mathcal{H}_k}
     {C_{\text{cc}}}{k = 0 \land c \geq 0}{c + N\mathcal{H}_k},
\]
certifying $\mathbb{E}[c_{\text{final}}] \geq c_0 + N\mathcal{H}_{k_0}$;
from a fresh start, $\mathbb{E}[\text{\#draws}] \geq N \mathcal{H}_N$ ---
the classical coupon-collector bound, again exact. The certified rate
$1-\sfrac{1}{2N}$ is within a factor of two (in the exponent) of the true
rate $1-\sfrac1N$.

% ----------------------------------------------------------------------
\subsection{Escaping Spline}\label{app:bench-spline}
% ----------------------------------------------------------------------

\paragraph{Program.}
The \emph{escaping spline} (Hark et al.\ 2020), with $x$ ranging over the
naturals:
\[
  C_{\text{sp}} \;\defsym\;
  \WHILE{x > 0}{\;
    (x \mathrel{:=} 0) \oplus_{1/(x+1)} (x \mathrel{:=} x+1)
  \;}
\]
From $x \geq 1$ the loop exits in one step with probability
$\sfrac{1}{x+1}$ and otherwise drifts upward, making the exit ever less
likely. The tail probabilities telescope in closed form:
\[
  q_n(x)
  = \prod_{i=0}^{n-1} \frac{x+i}{x+i+1}
  = \frac{x}{x+n},
\]
so the loop is almost surely terminating, yet
$\mathbb{E}[T_x] = \sum_{n \geq 0} \tfrac{x}{x+n} = \infty$: almost-sure
termination with infinite expected runtime, the classical hard regime for
lower bounds. (It also illustrates the finiteness caveat of
\Cref{app:completeness}: for the counter specification the exact
pre-expectation $H$ is infinite.)

\paragraph{Specification.}
Take $\pred \equiv x \geq 0$, $\bexpr \equiv x > 0$, and $\ex \defsym 1$.
The target triple
\[
  \HT[lower]{x \geq 0}{1}{C_{\text{sp}}}{x = 0}{1}
\]
asserts $\WP{C_{\text{sp}}}{[x=0]\cdot 1} \geq 1$, i.e.\ almost-sure
termination.

\paragraph{\textsc{While-Mult} is provably inapplicable.}
By \Cref{prop:geo-tail}, any instantiation would give
$q_n(x) \leq \beta^n V(x)$ with $\beta \lt 1$ and $V(x) \lt \infty$. But
$q_n(x) = \tfrac{x}{x+n} \geq \tfrac{x}{2n}$ for $n \geq x$, while
$\beta^n V(x) \to 0$ geometrically --- a contradiction for $n$ large. The
polynomial tail forces the general rule.

\paragraph{Witnesses.}
The canonical witness sequence of \Cref{app:canonical-examples} is available
in closed form: take
\[
  r_n(x) \;\defsym\; [x \geq 1]\cdot\frac{x}{x+n}.
\]

\paragraph{Premise 1} (with equality; both successors satisfy $\pred$):
\[
  \WP{C_{\text{body}}}{[\pred]\cdot 1}(x)
  = \frac{1}{x+1}\cdot 1 + \frac{x}{x+1}\cdot 1
  = 1
  \;\geq\; 1.
  \checkmark
\]

\paragraph{Premise 2.}
On $\bexpr\land\pred$, i.e.\ $x \geq 1$:
$r_0(x) = \tfrac{x}{x} = 1 \geq \ex$. $\checkmark$

\paragraph{Premise 3} (with equality): for $x \geq 1$, using $r_n(0) = 0$,
\[
  \WP{C_{\text{body}}}{r_n}(x)
  = \frac{1}{x+1}\cdot r_n(0) + \frac{x}{x+1}\cdot r_n(x+1)
  = \frac{x}{x+1}\cdot\frac{x+1}{x+n+1}
  = \frac{x}{x+n+1}
  = r_{n+1}(x).
  \checkmark
\]

\paragraph{Premise 4.}
For each fixed $x \geq 1$:
$\lim_{n\to\infty} \tfrac{x}{x+n} = 0$. $\checkmark$

\paragraph{Conclusion.}
$\vdash\; \HT[lower]{x \geq 0}{1}{C_{\text{sp}}}{x = 0}{1}$: the loop
terminates with probability $1$, certified entirely within the logic by
elementary algebra --- every premise is discharged by a one-line computation,
with no auxiliary termination argument. This benchmark exercises the general
rule at full strength precisely where the derived rule provably fails.

% ----------------------------------------------------------------------
\subsection{Unbounded Symmetric Random Walk}\label{app:bench-rw}
% ----------------------------------------------------------------------

\paragraph{Program.}
\[
  C_{\text{rw}} \;\defsym\;
  \WHILE{x > 0}{\;
    (x \mathrel{:=} x-1) \oplus_{1/2} (x \mathrel{:=} x+1)
  \;}
\]
The motivating example of the lower-bound literature: the walk is almost
surely terminating (recurrence of the one-dimensional symmetric random
walk), yet $\mathbb{E}[T_x] = \infty$ and the tail is only polynomial. By
the classical ballot-problem identity, $q_{2m}(1) = \binom{2m}{m}4^{-m}$,
and with the standard central-binomial bound
$\binom{2m}{m} \geq 4^m/(2\sqrt{m})$ (together with $q_{2m-1}(1) =
q_{2m}(1)$, as $T_1$ only takes odd values) one gets
\[
  q_n(1) \;\geq\; \frac{1}{2\sqrt{n}}
  \qquad (n \geq 1).
\]
Hence, by \Cref{prop:geo-tail}, \textsc{While-Mult} is provably inapplicable
to the termination triple, exactly as for the escaping spline.

\paragraph{Specification.}
$\pred \equiv x \geq 0$, $\bexpr \equiv x > 0$, $\ex \defsym 1$, and the
target triple $\HT[lower]{x \geq 0}{1}{C_{\text{rw}}}{x = 0}{1}$
(almost-sure termination).

\paragraph{Witnesses.}
Unlike the spline, the canonical witnesses have no elementary closed form.
We define them by the one-step recursion
\[
  r_0(x) \defsym [x \geq 1],
  \qquad
  r_{n+1}(x) \defsym [x \geq 1]\cdot
    \Bigl(\tfrac12\, r_n(x-1) + \tfrac12\, r_n(x+1)\Bigr),
\]
so that $r_n(x) = q_n(x) = \Pr(T_x > n)$ by induction on $n$.

\paragraph{Premises 1--3.}
Premise~1: $\WP{C_{\text{body}}}{[\pred]\cdot 1} = 1 \geq 1$ on
$\bexpr\land\pred$, as both successors satisfy $\pred$. Premise~2:
$r_0 = 1 \geq \ex$ on $x \geq 1$. Premise~3 holds \emph{with equality by
construction}: for $x \geq 1$,
$\WP{C_{\text{body}}}{r_n}(x) = \tfrac12 r_n(x-1) + \tfrac12 r_n(x+1)
 = r_{n+1}(x)$. $\checkmark$

\paragraph{Premise 4.}
For fixed $x \geq 1$, $\;\lim_{n\to\infty} r_n(x) = \Pr(T_x = \infty)$, so
Premise~4 is \emph{equivalent} to almost-sure termination of the walk. This
is true, but it is a genuinely external fact: it must be imported, either
from the classical recurrence argument or from a dedicated AST rule (e.g.\
the supermartingale variant rule of McIver et al.\ 2018, or the proof rules
of Majumdar and Sathiyanarayana 2025).

\paragraph{Discussion.}
This benchmark marks the honest boundary of the method. The canonical
witnesses exist, Premises~1--3 hold by construction, and the entire
difficulty of the example is concentrated in Premise~4: the rule reduces the
quantitative statement ``the termination probability is at least $1$'' to
the qualitative statement ``the loop terminates almost surely'', and no
further. This is the relative-completeness phenomenon in concrete form ---
any proof system for lower bounds must certify almost-sure termination
somewhere, and the rule makes that dependency explicit instead of hiding it.
Contrast with \Cref{app:bench-spline}, where the same recipe succeeded
outright because $q_n$ admitted a closed form with a directly computable
limit.


\end{document}
